\documentclass[10pt,a4paper]{article}
\usepackage[margin=1.75cm]{geometry}
\usepackage{amsmath,amssymb,booktabs,longtable,array,multirow}
\usepackage{graphicx,xcolor,enumitem,listings,caption,subcaption,needspace,etoolbox}
\usepackage{tikz}
\usetikzlibrary{arrows.meta,positioning,fit,shapes.geometric,backgrounds,calc}
\usepackage{pdflscape}
\usepackage[colorlinks=true,linkcolor=blue!55!black,citecolor=green!40!black,urlcolor=blue!65!black]{hyperref}
\usepackage[capitalize,nameinlink]{cleveref}
\usepackage{microtype}
\usepackage{fancyhdr}
\usepackage{titlesec}
\usepackage{natbib}
\usepackage{doi}
\usepackage[T1]{fontenc}
\usepackage{lmodern}
\setlength{\parskip}{0.35em}
\setlength{\parindent}{0pt}
\setlist{nosep,leftmargin=*}
\pagestyle{fancy}\fancyhf{}\lhead{Long Context, Retrieval, and Structured Access}\rhead{PRAttention-centered tutorial}\cfoot{\thepage}
\definecolor{codebg}{RGB}{246,248,250}
\definecolor{localgray}{RGB}{105,112,120}
\definecolor{retrievalblue}{RGB}{48,105,190}
\definecolor{memorygreen}{RGB}{45,145,88}
\definecolor{structurepurple}{RGB}{126,77,166}
\definecolor{agentred}{RGB}{190,63,63}
\definecolor{praorange}{RGB}{224,126,38}
\lstset{basicstyle=\ttfamily\scriptsize,backgroundcolor=\color{codebg},frame=single,breaklines=true,columns=fullflexible,showstringspaces=false,language=Python}
\BeforeBeginEnvironment{lstlisting}{\Needspace{15\baselineskip}}
\newcommand{\methodflow}[6]{%
\par\nopagebreak[4]\begin{center}
\begin{tikzpicture}[node distance=3mm and 5mm,
  stage/.style={draw=#6!80!black,fill=#6!8,rounded corners,align=center,
    inner xsep=1.5pt,minimum height=7mm,text width=.155\linewidth,font=\scriptsize},
  flow/.style={-{Latex[length=1.7mm]},draw=#6!80!black,thick}]
\node[stage] (a) {#1};
\node[stage,right=of a] (b) {#2};
\node[stage,right=of b] (c) {#3};
\node[stage,right=of c] (d) {#4};
\node[stage,right=of d] (e) {#5};
\draw[flow] (a)--(b); \draw[flow] (b)--(c); \draw[flow] (c)--(d); \draw[flow] (d)--(e);
\end{tikzpicture}
\end{center}\par}
\title{\textbf{Long Context, Retrieval, Structured Indexing, and Search Agents:}\\An Educational Architecture Review Centered on Progressive Retrieval Attention}
\author{Jorge Sim\~ao\\EInnovator R\&D -- Center for the Sciences of Computation, Cognition, and Complexity\\\small Tutorial and research-design draft}
\date{August 2026}
\begin{document}
\maketitle
\begin{abstract}
Language models can gain access to more information by making the token window larger, restricting attention, carrying recurrent state, consulting an external memory, retrieving text before generation, retrieving latent key/value representations during attention, navigating a structured document index, or invoking search tools iteratively. These mechanisms are often discussed as separate subfields even though they answer the same architectural questions: \emph{what is stored, how is it named, when is relevance decided, where is retrieved information materialized, and how can its causal contribution be tested?}

This tutorial organizes the field into seven categories and uses Progressive Retrieval Attention (PRAttention, PRA) as a unifying reference point. PRA represents external context through explicit URI-like handles, resolves selected objects into layer-specific key/value memory, and permits demand-driven access through a dedicated attention branch. We first reconstruct the PRA position and standalone prototype papers, including their preliminary WikiText-2 results and counterfactual evaluation doctrine. We then explain representative approaches using motivation, formalism, pseudocode, PyTorch sketches, diagrams, original experiments, research lineage, downstream influence, relation to PRA, and concrete comparison experiments. The final sections develop a common design space, propose fair experimental matrices, discuss provenance and the 2026 PaperTrail interface, and derive extensions that combine hierarchical references, learned routing, cache reuse, and agentic escalation. The intended audience is students, researchers, and implementers who need both a conceptual map and reusable experimental designs.
\end{abstract}
\tableofcontents
\newpage

\section{Introduction}
The canonical Transformer converts a sequence into contextual representations using dense self-attention. Its flexibility comes with a quadratic interaction matrix and an implicit assumption: all potentially useful information is already present as a flat token sequence. Long-document and knowledge-intensive tasks violate that assumption in several distinct ways. The information may be too large to fit; relevance may be sparse; the useful evidence may be hierarchically organized; a query may evolve during reasoning; or the relevant source may not be known until generation has begun.

The literature responds with mechanisms at different system boundaries. Sparse attention changes the graph within a prompt. Recurrence carries state across prompt segments. Differentiable memory compresses or stores prior activations. Retrieval-augmented generation chooses passages and inserts them as text. Retrieval-attention systems index hidden states or key/value vectors. Structural indexes let a model navigate sections and summaries. Search agents act outside the model through repeated queries and tool calls. PRA is best understood not as replacing all of these, but as proposing a different interface between them: lightweight references remain in the local stream; named external objects are resolved only when useful; and the materialized representation is a layer-local latent memory rather than necessarily prompt text.

\subsection{Seven categories}
\begin{enumerate}
\item \textbf{Dense, sparse, and approximate long-context attention}: extend the visible token graph.
\item \textbf{Recurrence and compression}: carry or summarize state across segments.
\item \textbf{Differentiable external memory}: read and write learned memory slots or episodic stores.
\item \textbf{Retrieval-augmented pretraining and generation}: retrieve textual evidence before or during generation.
\item \textbf{Attention-level and KV retrieval}: retrieve hidden states or key/value memories directly into attention.
\item \textbf{Structural and hierarchical indexing}: navigate document, graph, or summary structure rather than flat chunks.
\item \textbf{Agents with search tools}: iteratively formulate queries, inspect results, follow links, and decide when to stop.
\end{enumerate}

\subsection{A shared vocabulary}
Let $x_{1:T}$ be local tokens, $q_t$ a query representation at generation step $t$, $\mathcal{S}$ an external store, $R(q_t,\mathcal{S})$ a selection operator, and $M_t$ the materialized memory. A general access architecture is
\begin{align}
C_t &= R(q_t,\mathcal{S};\pi),\\
M_t &= \operatorname{materialize}(C_t;\rho),\\
h'_t &= \operatorname{combine}(h_t,M_t;\gamma),
\end{align}
where $\pi$ is a selection policy, $\rho$ determines representation (tokens, summaries, hidden states, KV tensors, tool observations), and $\gamma$ determines fusion. Every method in this review can be located by its choices for $\mathcal{S},R,\rho,$ and $\gamma$.

\section{Progressive Retrieval Attention}
\subsection{PRA as a Worked Decomposition}

Progressive Retrieval Attention (PRAttention, PRA) begins with a separation that ordinary
prompt construction leaves implicit: information may be \emph{logically available} without
participating in every model operation. A prompt can advertise a document, prior prompt
prefix, code object, or record through a stable handle. The runtime can retain a compact
search view for many such objects while admitting full-detail memory for only a selected
working set. The model therefore sees two context sizes:
\begin{equation}
N_{\rm logical}=\left|\bigcup_{u\in\mathcal U}z_u\right|,
\qquad
N_{\rm active}(t,\ell)=T_{\rm local}+\sum_{c\in A_{t,\ell}}|c|,
\end{equation}
where $z_u$ is the source bound to URI $u$, $A_{t,\ell}$ is the set of chunks active at
step $t$ and layer $\ell$, and $T_{\rm local}$ is the direct sequence length. Bounded
execution requires a configured upper bound on $N_{\rm active}$, not a promise that the
logical store is small.

This is an access contract with five operations:
\begin{enumerate}
\item \textbf{name}: a handle identifies an object or a request-local prompt prefix;
\item \textbf{resolve}: a policy-aware resolver binds identity and version to content;
\item \textbf{index}: compact gists make chunks searchable without treating the gist as
the evidence itself;
\item \textbf{materialize}: selected chunks expose full-detail model-native K/V; and
\item \textbf{consume}: selected native K/V enters chosen pretrained attention layers.
\end{enumerate}
PRA's novelty claim is this decomposition and its bounded composition, not the invention of
top-$k$, K/V caching, cross-attention, or retrieval. The position, standalone, and pretrained
studies evaluate different pieces of the contract \citep{simao2026praArchitecture,simao2026praStandalone,
simao2026praPosition,simao2026praHF}. They do not establish a universal advantage over long
context or RAG.

\methodflow{logical objects + local tail}{model-safe native encoding}{compact semantic index}{budgeted native K/V}{selected pretrained layers}{praorange}

\subsection{Identity, Resolution, and Bounded Logical Memory}

Let the context graph be $\mathcal G=(\mathcal V,\mathcal E)$. A node
$v=(u,z,m,R,\nu)$ contains URI $u$, content $z$, metadata $m$, an explicit child-reference
table $R$, and version $\nu$. An edge exists because a reference table names a child, not
because two URI strings share a prefix. A resolver remains outside the neural layer:
\begin{equation}
\rho(u,p,t)\longrightarrow(z,m,R,\nu),
\end{equation}
where $p$ is an authorization domain and $t$ is logical time. The separation permits files,
database rows, prompt history, and document spans to share an interface without pretending
that they share a physical store.

A cache key for reusable neural payload must bind at least URI, content version, model,
consumer layer, encoding policy, positional provenance, and authorization domain. Otherwise
identical names can alias different content or a K/V tensor can be consumed under a model or
coordinate system that did not create it. Request-private caches avoid cross-request leakage
but do not demonstrate shared reuse. Persistent reuse therefore adds invalidation, tenancy,
and provenance obligations beyond the attention mechanism.

Long prompts use the same contract without requiring an explicit document URI. The recent
tail remains in direct causal attention; an older prefix becomes a request-local implicit
\texttt{\_\_head\_\_} object with its own chunk, routing, and budget parameters. This makes
old prompt information recoverable without keeping all of it active. It does not extend the
backbone's learned positional range, and generated-token rollover is a separate streaming
policy.

Recursive graph construction is bounded by shared object, depth, branch, source-token, and
materialization budgets. Child-first construction allows a parent representation to use
completed child memory where configured, while cycle states prevent unbounded traversal.
These mechanics establish safe graph handling. A learned policy that follows evidence from
one retrieved object to another has not yet been established.

\subsection{Encoding Granularity, Routing Granularity, and Native Payload}

One source partition cannot optimize every stage. An \emph{encoding block} must fit the
backbone's model-safe context and preserve enough neighboring text for useful hidden states.
A smaller \emph{routing chunk} supplies search resolution. Several routing chunks may point
into one contiguous encoded block, and their selected spans recover the corresponding native
keys and values. Thus independent 32-token encoding and 32-token routing is not equivalent to
contiguous 512-token encoding followed by 32-token routing.

For layer $\ell$, source block hidden states $H^{(\ell)}_c$ are projected through the
pretrained attention module:
\begin{align}
K^{(\ell)}_c &= \operatorname{RoPE}_{p_c}
  \left(H^{(\ell)}_c W^{(\ell)}_K\right), &
V^{(\ell)}_c &= H^{(\ell)}_c W^{(\ell)}_V.
\end{align}
These tensors retain the native K/V head layout. The runtime does not insert a separate
memory-value projection in the default path. After selection, it concatenates admitted
native payload with local native K/V and invokes the model's ordinary eager attention
function with a mask that preserves local causality and permits access to memory:
\begin{align}
\widetilde K^{(\ell)}&=[K^{(\ell)}_{A};K^{(\ell)}_{\rm local}], &
\widetilde V^{(\ell)}&=[V^{(\ell)}_{A};V^{(\ell)}_{\rm local}],\\
O^{(\ell)}&=\operatorname{NativeAttn}
  \left(Q^{(\ell)}_{\rm local},\widetilde K^{(\ell)},
  \widetilde V^{(\ell)},\widetilde M\right).
\end{align}
The disabled branch delegates to the unmodified attention module. Exact-bypass tests are
therefore a compatibility gate: no selected memory must yield the original logits, not an
approximation to them. For fixed selected identities, the tested K, V, and native attention
output are bit-exact in the positional study \citep{simao2026praPosition}.

This native path differs from the historical standalone cross-attention variant. That
variant encoded reference K/V through a distinct memory branch and output projection; it is
useful as an adaptation control but breaks native-payload equivalence. The survey treats it
as a secondary design point rather than PRA's default transport.

\subsection{Gists and Search Geometry}

Routing should inspect a representation much smaller than full payload. For chunk $c$, a
simple gist partitions hidden states into $G$ contiguous subchunks and mean-pools each:
\begin{equation}
g_{c,j}=\frac{1}{|S_{c,j}|}\sum_{i\in S_{c,j}}h_i,
\qquad j\in\{1,\ldots,G\}.
\end{equation}
A query $q$ scores each chunk by maximum, mean, or log-sum-exp aggregation over its gists.
Reference-first routing can rank URI-level gists before chunk gists. Fixed-$k$ bounds the
absolute payload; fixed-fraction selection holds sparsity pressure approximately constant
as the candidate count grows. Both should be reported because a constant $k$ becomes a
harder retrieval task at larger scales.

The important lesson from positional and pretrained experiments is that a native attention
key is not automatically a good semantic index. Centered multi-gists increasingly reproduce
native token-Q/K ranking, yet their semantic retrieval quality remains almost flat. A learned
32-dimensional asymmetric metric shows the reverse: strong semantic ranking with low native
Q/K fidelity. The robust decomposition is therefore
\begin{equation}
\underbrace{r_\phi(q,g_c)}_{\text{semantic discovery}}
\quad\longrightarrow\quad
\underbrace{(K_c,V_c)}_{\text{native execution}}.
\end{equation}
The router may use learned semantic geometry while the selected payload remains untouched
native K/V. This separation also prevents a lossy routing summary from silently becoming the
memory content.

\subsection{Position Is Not Context}

Independently encoded chunks introduce two different errors. \emph{Positional reset} assigns
each chunk positions beginning at zero. A source-relative logical offset repairs this error
for learned absolute, sinusoidal, and RoPE coordinates. For RoPE, cached pre-rotation keys can
also be rebound algebraically to a new offset, although deferred rotation has a measured
runtime cost.

\emph{Contextual fragmentation} remains after coordinates are repaired. A token encoded
without its left source context has a different hidden state from the same token inside the
contiguous document. Overlap or larger grouped encoding blocks reduce this mismatch at added
encoding and storage cost. The resulting rule is simple: preserve source-relative coordinates,
but do not describe coordinate correctness as recovered contextualization.

The measured geometry results make the distinction concrete. Source-relative offsets remove
layer-zero reset error in all paired tests and improve deeper representations. Increasing
centered-RoPE gists from one to eight raises native Q/K fidelity from $.794$ to $.963$, while
semantic sparse AUC remains approximately $.651$. A learned router reaches $.971$ semantic
AUC at only $.178$ Q/K fidelity. Position, contextual representation, semantic search, and
native attention are four related but distinct objects.

\subsection{Depth, Integration, Decision, and Behavior}

Correct routing supplies evidence; it does not guarantee that a pretrained decoder knows how
to use that evidence. PRA can be enabled at selected layers $\mathcal L_{\rm PRA}$. Pretrained
experiments find a non-monotone depth effect: a late layer band can improve memory use, while
injection into every layer can be disruptive. This turns layer placement into an integration
variable rather than an implementation detail.

An oracle intervention replaces only selected identities while preserving transport, budget,
and consumer. If oracle memory helps while routed memory does not, discovery is the bottleneck.
If oracle memory also fails, representation or integration is implicated. QASPER provides a
measured instance of this ladder: answer-polarity success progresses from 25.0\% under frozen
routed PRA to 62.5\% under frozen oracle memory and $72.5\pm5.6$\% with a small residual
adapter. Blind behavioral judges also favor the residual-adapted QASPER outputs. The result
shows that one-shot memory can help and that consumer calibration still matters.

The same experiment also shows why likelihood is insufficient. A model can improve
teacher-forced loss without completing an answer, and a token-level polarity improvement can
coexist with truncation. Evaluation therefore needs answer-critical logit margins, completion
rate, lexical metrics, task semantics, and blinded behavioral comparison. On HotpotQA, outputs
often approach a relevant entity or relation but miss another required hop or endpoint. This
is evidence of a one-shot relational boundary, not proof that more decoder adaptation alone
will solve multi-hop discovery.

\subsection{Evidence Ledger}

Table~\ref{tab:pra-evidence-ledger} summarizes the current empirical boundary. Results come
from different model scales and protocols and should not be pooled as one benchmark score.
Their role in this survey is to identify which links in the access chain have been tested.

\begin{table}[htbp]
\centering\scriptsize
\caption{PRA evidence by computational boundary. ``Measured'' does not imply production
validation or universal superiority.}
\label{tab:pra-evidence-ledger}
\begin{tabular}{>{\raggedright\arraybackslash}p{2.5cm}>{\raggedright\arraybackslash}p{6.1cm}>{\raggedright\arraybackslash}p{5.8cm}}
\toprule
Boundary & Measured result & Interpretation and limit \\\midrule
Ordinary path & Disabled adapters preserve the host path exactly across standalone and
pretrained-family tests & Compatibility gate; it says nothing about memory utility \\
Position & Source-relative offsets repair reset fragmentation; native pre/post/deferred RoPE
paths agree under matched coordinates & Coordinate correctness is established; contextual
fragmentation remains \\
Logical versus active & At 64--256 addressable units, fixed $k=8$ holds physical selected K/V
near 11--13 tokens and reduces active fraction to 4.4--4.9\% at 256 & Bounded native transport
on controlled probes, not unrestricted long-document QA \\
Routing scale & At about 12.5\% selection, target coverage remains .756--.772 on
HotpotQA-derived and .794--.803 on QASPER-derived probes through 256 units & Stable fixed-fraction
quality; substantial router headroom remains \\
Search geometry & Centered $G=1\rightarrow8$ improves Q/K fidelity $.794\rightarrow.963$ with
flat semantic AUC; learned 32-D routing reaches .971 AUC & Semantic search should be separated
from native execution \\
Pretrained retrofit & Qwen, Llama-family, and Gemma adapters preserve disabled semantics;
selected native K/V uses family-native head layouts & Cross-family mechanical portability at
roughly one-billion-parameter scale, not production coverage \\
Depth and use & Late-band memory is useful; indiscriminate all-layer activation is harmful in
the measured setting & Integration is depth-sensitive \\
One-shot QASPER & Routed frozen 25.0\%, oracle frozen 62.5\%, residual adapter
$72.5\pm5.6$\% polarity; behavioral judges favor adaptation & Positive one-shot regime on a
small fixed cohort \\
One-shot HotpotQA & Relation-near outputs frequently miss another required evidence item or
endpoint & Motivates iterative closure; does not validate it \\
Serving & Packed tensor routing uses vectorized scoring and \texttt{topk}; 256-way warm
selection is 5.5--11.6 ms in the standalone prototype & Better than exhaustive scalar routing,
but no optimized serving-kernel claim \\
\bottomrule
\end{tabular}
\end{table}

\subsection{Systems Accounting}

Let $G$ be the number of routing gists, $M$ selected native K/V tokens, $L_p$ PRA-enabled
layers, and $w$ bytes per scalar. Ignoring head constants, exact dense gist scoring costs
$O(Gd)$ and active memory attention costs $O(L_pT_{\rm local}Md)$. Resident selected payload
costs approximately
\begin{equation}
B_{\rm active}\simeq 2wL_pMd,
\end{equation}
while the backing cache may be larger and may reside on CPU or another storage tier. Reducing
active attention does not automatically reduce total storage. Likewise, a short prompt is not
an efficiency gain if every source is cold-encoded for every request.

A systems report must separate source resolution, tokenization, model-safe encoding, index
construction, routing, host-to-device transfer, native attention, synchronization, and
generation. It should report cold, warm, and authorized shared-cache regimes; logical tokens;
available and selected gists; physical selected K/V; peak accelerator and host bytes; time to
first token; inter-token latency; and avoided encoding work. Packed \texttt{torch.topk} is a
prototype optimization, not evidence of an asymptotically indexed router or fused sparse
attention kernel.

\subsection{Causal Evaluation Doctrine}

For loss $L$, define
\begin{equation}
\Delta_{\rm use}=L_{\rm disabled}-L_{\rm valid},\qquad
\Delta_{\rm content}=L_{\rm shuffled}-L_{\rm valid}.
\end{equation}
A positive $\Delta_{\rm use}$ only shows that the active branch changes the model usefully.
Content-causal retrieval additionally requires valid evidence to outperform shuffled,
irrelevant, stale, and matched-length controls. Oracle identity measures the upper bound of
the same consumer; direct text measures whether the backbone can solve the item when evidence
is local. All-reference and gist-only conditions isolate routing and detail bottlenecks.

These controls should be crossed with routing budget, candidate count, layer band, and
materialization mode. Report reference recall, all-required-evidence recall, physical K/V
fraction, answer-critical logit margins, task output, and free-running behavior. Attention
weights and lexical overlap are diagnostics, not causal proof.

\subsection{Architecture Ladder: Established, Next, and Later}

The PRA research program now has three distinct levels.

\paragraph{Established in the current evidence.}
Bounded logical memory, model-safe source encoding, separate routing chunks, semantic gist
routing, full-detail native-KV payload, pretrained retrofit with exact disabled paths, and
depth-sensitive integration are implemented and measured. The evidence is strongest for
mechanical fidelity, controlled scaling, routing diagnostics, and a small one-shot QASPER
regime.

\paragraph{Next: associative closure over gists.}
One-shot routing computes $C_1=R(q_0,\mathcal I)$. A proposed bounded iterative policy lets
retrieved gist content update the query and expose another relation:
\begin{align}
C_h &= R(q_{h-1},\mathcal I\setminus V_{h-1}),\\
q_h &= U(q_{h-1},G(C_h)),\qquad h=1,\ldots,H,\\
A &= \operatorname{BudgetSelect}\!\left(\bigcup_{h=1}^{H}C_h,B\right).
\end{align}
Here $V_{h-1}$ prevents cycles, $H$ bounds hops, and $B$ fixes the final unique native-KV
budget. The retrieved objects form an evidence path rather than an unordered bag. HotpotQA is
the natural primary test because the current failure is often a missing relation or endpoint;
QASPER is a one-shot control. This is \emph{associative closure}: a retrieved concept helps
discover another source. It is not yet an empirical PRA result and should not be described as
full neural reasoning.

\paragraph{Later: progressive and smart materialization.}
After relevant chunks are found, a second policy may decide which native detail to admit.
Partial pages, progressive escalation, token-level bounds, and head/layer-specific budgets are
possible mechanisms. Under the \textbf{redundancy hypothesis}, less detail yields equivalent
quality because much selected K/V is unnecessary. Under the \textbf{overload or dilution
hypothesis}, less detail improves quality because irrelevant detail competes with causal
evidence, a concern consistent with positional sensitivity in long contexts
\citep{liu2024lostmiddle}.
If active K/V falls while quality is unchanged, the first explanation is supported. If quality
improves, the second becomes plausible but still requires content-removal controls. If quality
falls, omitted detail was useful or the allocator was wrong. These are Paper~3 hypotheses,
not consequences of the one-shot results.

The ladder preserves the central thesis without overextending it: PRA is an emerging
decomposition of bounded memory computation. It has established a native one-shot working-set
mechanism; it has not yet established relational closure, optimal materialization, or a
production serving advantage.

The next seven sections apply the access-contract vocabulary to representative systems.
To keep the mathematics comparable, $T$ denotes sequence length, $d$ model width, and
$d_h$ the width of one attention head. The matrices $Q$, $K$, and $V$ contain projected
queries, keys, and values; $M_{\rm causal}$ masks positions that a causal decoder may not
read. Method-specific symbols are introduced where they appear. Complexity expressions
describe the dominant conceptual operation and omit constant head and implementation
factors unless those factors are themselves the subject of the method.

\section{Category I: Dense, Sparse, and Approximate Long-Context Attention}

This family changes the token-interaction graph after information has already been
materialized as one sequence.  It is therefore the cleanest ``all evidence local'' control
for PRA: selection is either absent or encoded in the attention mask, object identity is
reduced to position, and no cross-request materialization is implied.

\subsection{Dense attention and positional extension}

\textbf{Motivation and mechanism.} Dense causal attention is the least structured but
most general access rule: every position can read every earlier position.  RoPE scaling,
position interpolation, and long-context continued pretraining extend the usable position
range without changing this graph.  For $T$ tokens,
\begin{equation}
A=\operatorname{softmax}((QK^\top+M_{\rm causal})/\sqrt{d_h})V,
\qquad \text{time }O(T^2d),\;\text{attention storage }O(T^2).
\end{equation}
Positional extrapolation and attention efficiency are different problems. A model may
accept a long index while still allocating quadratic
work, or execute efficiently while failing to use distant positions.

\begin{lstlisting}[caption={Dense causal attention; PyTorch uses optimized kernels when possible}]
def dense_causal(q, k, v):              # [B,H,T,D]
    return F.scaled_dot_product_attention(q, k, v, is_causal=True)
\end{lstlisting}
\methodflow{all source tokens}{tokenize eagerly}{dense causal graph}{layer-by-layer KV}{decoder output}{localgray}

\textbf{Evidence, failure modes, and lineage.} Standard language modeling and long-context
adaptations use perplexity, long-document QA, retrieval probes, and throughput/memory
sweeps.  Dense attention is a strong baseline when sources are short or densely relevant,
and modern fused kernels make it unexpectedly competitive at moderate lengths.  Its
failures include quadratic cost, distraction, positional degradation, and lack of stable
object identity or reuse beyond exact prefixes.  Needle tests diagnose reachability but do
not establish synthesis or citation behavior.

\textbf{Relation and comparison to PRA.} Dense context performs eager materialization;
PRA advertises a potentially larger logical context and materializes selected objects.
Compare them on the same source corpus under (i) equal visible-plus-retrieved tokens,
(ii) equal FLOPs, (iii) equal peak KV bytes, and (iv) equal synchronized latency.  Sweep
evidence density and object reuse.  Dense attention should be expected to win when most
tokens matter once; PRA's hypothesis is strongest when few named objects serve many
queries.  The all-content-local condition is also PRA's transport upper control.

\subsection{Longformer}

\textbf{Motivation and mechanism.} Longformer observes that document dependencies are
mostly local but a few task positions need global connectivity.  A token attends within a
window and, optionally, to a set $\mathcal G$ of global positions
\citep{beltagy2020longformer}:
\begin{equation}
\mathcal N(i)=\{j:|i-j|\leq w/2\}\cup\mathcal G,
\qquad O(Tw+T|\mathcal G|).
\end{equation}
Local layers propagate information gradually; global tokens create short paths across the
document and can be task-selected (for example, question tokens in QA).  Unlike truncation,
the model retains an explicit sparse graph over the entire input.

\begin{lstlisting}[caption={Educational Longformer mask; production kernels avoid materializing $T^2$}]
def longformer_mask(T, window, global_pos, device):
    i = torch.arange(T, device=device)
    local = (i[:, None] - i[None, :]).abs() <= window // 2
    g = torch.zeros(T, dtype=torch.bool, device=device)
    g[global_pos] = True
    return local | g[:, None] | g[None, :]
\end{lstlisting}
\methodflow{long token sequence}{sliding window}{task global tokens}{sparse self-attention}{document state}{localgray}

\textbf{Evidence and paper trail.} The original paper studies character language modeling,
long-document classification, and QA; its encoder-decoder continuation LED applies the
pattern to long-input generation.  Comparisons include RoBERTa-style dense baselines and
task-specific long-document systems, using bits-per-character, accuracy/F1, and generation
metrics.  Longformer became a template for domain adaptations and local/global attention.
Its main failure is structural: useful long-range routes must be known through the global
mask or emerge through multiple local hops.  Fixed windows can split tables, citations,
or cross-section evidence, and global positions can become a bottleneck.

\textbf{PRA relation and matched experiment.} Longformer makes relevance a sparse edge
decision over already materialized tokens; PRA makes it an object-admission decision
before latent transport.  Give both the same documents and vary evidence sparsity,
distance, and whether section boundaries are known.  Match the number of attended key
positions: Longformer local/global edges versus PRA top-$k$ object tokens.  Add shuffled
global markers and shuffled URI content.  This separates sensitivity to graph placement
from sensitivity to object semantics.  A hybrid can use Longformer locally inside a large
resolved object and PRA between objects.

\subsection{BigBird}

\textbf{Motivation and mechanism.} BigBird combines window, random, and global edges so
that each token has bounded degree while the graph mixes rapidly
\citep{zaheer2020bigbird}:
\begin{equation}
\mathcal N(i)=\mathcal N_{\rm window}(i)\cup
\mathcal N_{\rm random}(i)\cup\mathcal G.
\end{equation}
The random component supplies shortcuts without content-dependent routing; global tokens
connect task anchors to all positions.  Under the paper's graph assumptions this pattern
retains strong expressivity results while reducing sequence complexity to linear in $T$
for fixed block sizes.

\begin{lstlisting}[caption={Block-level BigBird adjacency sketch}]
def bigbird_blocks(n_blocks, radius, random_k, global_blocks=(0,)):
    edges = []
    for i in range(n_blocks):
        local = range(max(0, i-radius), min(n_blocks, i+radius+1))
        random = torch.randperm(n_blocks)[:random_k].tolist()
        edges.append(sorted(set(local) | set(random) | set(global_blocks)))
    return edges
\end{lstlisting}
\methodflow{blocked sequence}{local + random edges}{global task blocks}{sparse attention}{mixed document state}{localgray}

\textbf{Evidence, ablations, and limitations.} BigBird evaluates natural-language QA and
summarization as well as genomic sequences, comparing dense BERT/RoBERTa-style models and
other efficient Transformers with task accuracy/F1, ROUGE, and sequence-scale behavior.
The meaningful ablation is the contribution of local, random, and global edges.  Random
connectivity is inexpensive on average but can miss a particular necessary route;
global tokens require a placement rule, and sparse kernels do not automatically make an
otherwise identical pretrained model work at new lengths.

\textbf{PRA relation and experiment.} BigBird constructs content-independent shortcuts
inside one anonymous token graph.  PRA constructs content-dependent edges to explicitly
named external objects.  Compare a matched sparse-edge budget in which random/global block
edges and top-$k$ URI objects expose the same number of keys.  Sweep distractor blocks,
wrong handles, evidence distance, and repeated queries.  Report cold and warm materialization
for PRA; otherwise BigBird is unfairly charged only for attention while PRA hides encoding.
An instructive hybrid replaces random inter-document edges with URI-restricted routing and
retains BigBird within each selected document.

\subsection{Performer}

\textbf{Motivation and mechanism.} Performer approximates softmax attention with positive
random features (FAVOR+), exchanging a $T\times T$ matrix for associative summaries
\citep{choromanski2021performer}.  If
$\exp(q^\top k)\approx\phi(q)^\top\phi(k)$, then
\begin{equation}
\operatorname{Attn}(Q,K,V)\approx
D^{-1}\phi(Q)\bigl(\phi(K)^\top V\bigr),
\quad D=\operatorname{diag}(\phi(Q)(\phi(K)^\top\mathbf1)).
\end{equation}
For fixed feature count, time and memory are linear in sequence length.  Causal execution
uses prefix sums of $\phi(k_i)v_i^\top$ and $\phi(k_i)$.

\begin{lstlisting}[caption={Causal kernel-attention recurrence; phi must be a positive feature map}]
def causal_linear_attention(qf, kf, v, eps=1e-6):
    kv = torch.einsum('...tm,...td->...tmd', kf, v).cumsum(dim=-3)
    kz = kf.cumsum(dim=-2)
    num = torch.einsum('...tm,...tmd->...td', qf, kv)
    den = torch.einsum('...tm,...tm->...t', qf, kz).unsqueeze(-1)
    return num / den.clamp_min(eps)
\end{lstlisting}
\methodflow{all token Q/K/V}{random feature map}{associative KV statistic}{normalized query}{approximate output}{localgray}

\textbf{Evidence, failure modes, and lineage.} Performer reports image, text, and protein
experiments and measures both downstream quality and approximation/efficiency.  Its paper
trail lies in kernelized linear attention and random-feature estimation.  Feature variance,
finite-precision instability, causal implementation details, and inability to recover a
specific token after destructive aggregation are central failure modes.  Linear cost does
not mean selective relevance: every token contributes to the statistic.

\textbf{PRA relation and comparison.} Performer approximates computation after eager
materialization; PRA changes which objects enter computation.  The two are complementary:
use linear attention for the local stream or within a resolved memory.  A fair experiment
crosses exact versus FAVOR+ local attention with all-local versus PRA access, reports
approximation error and exact-value recall, and sweeps random-feature count under fixed
FLOPs.  PRA corruption tests remain necessary because a low approximation error says
nothing about whether the correct reference was selected.

\section{Category II: Recurrence and Compression}

Recurrence changes the temporal boundary: earlier segments leave behind state rather than
remaining as raw local tokens.  The central distinction from PRA is \emph{write-time
retention versus read-time dereferencing}.  Recurrent systems decide what history survives
as it passes; PRA assumes an addressable backing object can be fetched later.

\subsection{Transformer-XL}

\textbf{Motivation and mechanism.} Transformer-XL addresses fixed context and segment
fragmentation by carrying the previous segment's hidden states forward as stop-gradient
memory and using relative positional attention \citep{dai2019transformerxl}:
\begin{equation}
\widetilde H_{s}^{(\ell-1)}=[\operatorname{SG}(H_{s-1}^{(\ell-1)});H_s^{(\ell-1)}],
\quad Q=H_sW_Q,\quad K,V=\widetilde H_sW_{K,V}.
\end{equation}
Only current-segment states receive gradients through the recurrent boundary.  The memory
extends effective context without backpropagating across the entire history; relative
positions prevent a cached state from being confused with a current absolute position.

\begin{lstlisting}[caption={Segment recurrence; detach is the defining optimization boundary}]
def transformer_xl_segment(layer, x, old_memory, memory_len):
    kv_source = torch.cat([old_memory.detach(), x], dim=1)
    y = layer(query=x, key=kv_source, value=kv_source, relative_positions=True)
    new_memory = torch.cat([old_memory, x], dim=1)[:, -memory_len:].detach()
    return y, new_memory
\end{lstlisting}
\methodflow{current segment}{append cached history}{relative-position attention}{truncate memory}{next segment}{memorygreen}

\textbf{Evidence and limitations.} The paper evaluates WikiText-103, One Billion Word,
Penn Treebank, enwiki8, and text8 using perplexity or bits-per-character, alongside recurrent
and Transformer language-model baselines. It reports then-leading benchmark results and
substantially faster evaluation than recomputing overlapping windows. Ablations isolate
recurrence and relative positions.  Memory is chronological, anonymous, fixed-length, and
not revisable: evicted information is gone, irrelevant history consumes capacity, and a
detached boundary limits credit assignment.

\textbf{PRA relation and direct test.} Transformer-XL answers ``what recent state should
remain?''; PRA answers ``which named remote object should be loaded now?''  Compare them on
streams with either recency-correlated evidence or remote named evidence.  Match resident
KV bytes and total encoding FLOPs, sweep distance beyond the recurrent window, and include
repeated access to the same old object.  A hybrid keeps Transformer-XL memory for recent
discourse while letting a URI handle reload evicted or non-contiguous source material.

\subsection{Compressive Transformer}

\textbf{Motivation and mechanism.} Compressive Transformer extends Transformer-XL with a
second, lower-resolution memory.  States evicted from fine memory $M^m$ are transformed by
a compressor $C$ before entering compressed memory $M^c$ \citep{rae2020compressive}:
\begin{equation}
M^c_s=\operatorname{tail}_{n_c}([M^c_{s-1};C(E_{s-1})]),
\end{equation}
where $E_{s-1}$ denotes the fine states being evicted.  Attention reads both stores.  Pooling,
convolution, and usage-based compression trade temporal resolution for longer reach; auxiliary
reconstruction losses encourage retained predictive information.

\begin{lstlisting}[caption={Two-tier chronological memory update}]
def update_memory(fine, compressed, new_states, fine_cap, comp_cap, compress):
    merged = torch.cat([fine, new_states], dim=1)
    evicted, fine = merged[:, :-fine_cap], merged[:, -fine_cap:]
    compressed = torch.cat([compressed, compress(evicted)], dim=1)
    return fine.detach(), compressed[:, -comp_cap:].detach()
\end{lstlisting}
\methodflow{new segment}{fine recurrent memory}{compress evictions}{coarse long memory}{dual-memory attention}{memorygreen}

\textbf{Evidence and failure modes.} The original work studies long-range language modeling,
compares compression functions and auxiliary objectives, and uses perplexity plus attention
and memory diagnostics.  The method's honest trade-off is rate versus fidelity.  Pooling can
erase identifiers, superposition can merge distractors, learned compression can optimize
average prediction while losing rare facts, and chronological stores cannot directly
invalidate a stale semantic object.

\textbf{PRA relation and experiment.} PRA ordinarily reloads an original or cached object;
Compressive Transformer obligatorily summarizes passing history.  A useful hybrid attaches
several materializations to one URI: metadata, summary, compressed latent, and full KV, with
uncertainty-driven escalation.  Compare fine-only recurrence, dual-memory compression,
summary-only PRA, and full-object PRA under equal resident bytes.  Sweep compression ratio,
rare exact facts, semantic synthesis, and stale-object replacement.  Report whether escalation
recovers failures and charge full-object encoding in cold-cache trials.

\subsection{Recurrent Memory Transformer (RMT)}

\textbf{Motivation and mechanism.} RMT inserts a small set of learned memory tokens into
every segment.  Their output states become the memory-token inputs to the next segment,
creating an end-to-end learned bottleneck \citep{bulatov2022rmt}.  If $m_s$ are incoming
memory tokens and $x_s$ current tokens,
\begin{equation}
[\hat m_s,H_s,m_{s+1}]=\operatorname{Transformer}([m_s,x_s,m_s^{\rm write}]).
\end{equation}
The exact placement varies by implementation, but the conceptual operation is stable:
read old slots, process the segment, and write a fixed number of new slots.

\begin{lstlisting}[caption={RMT segment interface}]
def rmt_step(model, tokens, memory, write_tokens):
    sequence = torch.cat([memory, tokens, write_tokens], dim=1)
    hidden = model(sequence)
    next_memory = hidden[:, -write_tokens.size(1):]
    token_hidden = hidden[:, memory.size(1):-write_tokens.size(1)]
    return token_hidden, next_memory
\end{lstlisting}
\methodflow{memory tokens + segment}{full local processing}{write memory slots}{carry fixed bottleneck}{next segment}{memorygreen}

\textbf{Evidence, lineage, and failures.} RMT is evaluated on language modeling and
algorithmic/long-sequence tasks; later work demonstrates extrapolation on very long synthetic
problems.  Baselines include finite-window Transformers and recurrent-memory variants, with
loss, task accuracy, and length extrapolation as metrics.  The learned bottleneck is simple
and differentiable, but it must anticipate future needs at write time.  Rare details can be
overwritten, slot semantics lack stable provenance, and backpropagation depth or memory
curricula strongly affect training.

\textbf{PRA relation and experiment.} RMT learns \emph{what to remember while reading};
PRA learns \emph{what to fetch while answering}.  Compare both on matched memory bytes under
two datasets: one where future questions are predictable during ingestion and one where they
are not.  Substitute or reset RMT slots and shuffle PRA references to measure content use.
A hybrid stores an RMT summary in the URI metadata tier but preserves the source object for
read-time expansion when the summary is insufficient.

\subsection{Block-Recurrent Transformer}

\textbf{Motivation and mechanism.} Token-wise recurrence is sequential, whereas ordinary
self-attention is parallel but bounded by its active window. Block-Recurrent Transformer
places the recurrent boundary at a block of tokens. A set of state vectors $S_{b-1}$ and
the current token block $X_b$ exchange information through self- and cross-attention, then
LSTM-style gates form the state passed to the next block \citep{hutchins2022blockrecurrent}:
\begin{align}
\widetilde X_b&=\operatorname{TokenBlock}(X_b,S_{b-1}),\\
\widetilde S_b&=\operatorname{StateBlock}(S_{b-1},X_b),\\
S_b&=g_b\odot S_{b-1}+(1-g_b)\odot\widetilde S_b.
\end{align}
The exact cell contains several gates, but the computational principle is visible here:
parallel token processing occurs inside a block, while a fixed set of latent states carries
information between blocks. For fixed block and state sizes, total work grows linearly with
the number of processed tokens.

\begin{lstlisting}[caption={Simplified block-recurrent update}]
def block_recurrent_step(token_block, state, cell):
    token_out = cell.tokens_attend(token_block, state)
    state_candidate = cell.state_attends(state, token_block)
    gate = torch.sigmoid(cell.gate(torch.cat([state, state_candidate], -1)))
    next_state = gate * state + (1 - gate) * state_candidate
    return token_out, next_state
\end{lstlisting}
\methodflow{token block + state}{bidirectional block exchange}{gated state update}{fixed recurrent state}{next block}{memorygreen}

\textbf{Evidence and limitations.} The original study evaluates PG19 books, arXiv papers,
and GitHub source code against long-range Transformer-XL variants, reporting language-model
loss and execution behavior. The design uses accelerator-friendly block parallelism, but
the state remains an anonymous write-time bottleneck. A gate can preserve a feature without
preserving its source, version, or exact wording, and later queries cannot reload discarded
detail.

\textbf{PRA relation and experiment.} Block recurrence and PRA solve different halves of
the working-set problem. The recurrent state keeps a compact summary of the path already
read; PRA can dereference a named object that was never read or whose details no longer fit
the state. Compare recurrent state, PRA, and a hybrid under matched resident bytes on tasks
that independently vary temporal continuity and source addressability. State resets test
dependence on recurrent content; shuffled URI objects test dependence on PRA content.

\subsection{TransformerFAM}

\textbf{Motivation and mechanism.} A sliding-window transformer eventually loses access
to every old token, and a detached segment cache limits how gradients teach long-range
retention. TransformerFAM introduces a short Feedback Attention Memory (FAM) whose current
activations attend with each token block and whose updated activations feed the next block
\citep{hwang2024transformerfam}. If $F_{b-1}$ is the previous feedback state, one abstract
update is
\begin{equation}
[H_b,F_b]=\operatorname{MaskedAttn}([X_b,F_{b-1}],
                                    [X_b,M_b,F_{b-1}]),
\end{equation}
where $M_b$ denotes the ordinary block-sliding memory. The feedback queries are contextual:
they change after every block instead of using one static summary query. FAM adds no new
projection weights when it reuses the transformer's attention parameters, although it
changes training, masks, positions, and recurrent execution.

\begin{lstlisting}[caption={Feedback attention memory across blocks}]
def fam_step(block, sliding_memory, fam, attention, mask):
    queries = torch.cat([block, fam], dim=1)
    keys_values = torch.cat([block, sliding_memory, fam], dim=1)
    outputs = attention(queries, keys_values, keys_values, attn_mask=mask)
    return outputs[:, :block.size(1)], outputs[:, block.size(1):]
\end{lstlisting}
\methodflow{block + prior FAM}{joint masked attention}{contextual FAM update}{feedback to next block}{working-memory stream}{memorygreen}

\textbf{Evidence and limitations.} The paper studies models at 1B, 8B, and 24B parameters
on language modeling, passkey retrieval, NarrativeQA, and related long-context tasks. Its
ablations show that static summary queries and some simpler feedback variants do not retain
the passkey as well as contextual FAM. Improvements over block sliding attention can be
modest on several natural tasks, and the latent slots still lack stable object identity,
selective invalidation, or exact reconstruction after compression.

\textbf{PRA relation and experiment.} FAM is dynamic working memory; PRA is addressed
backing memory. A useful hybrid lets FAM carry the current reasoning state while PRA supplies
source detail. Compare a static summary token, contextual FAM, content-gist PRA, and the
hybrid. Use delayed exact recall, multi-hop synthesis, stale-source replacement, and unseen
question types. The key question is whether a small feedback state reduces PRA faults
without hiding content failures inside an uninterpretable latent bottleneck.

\subsection{Infini-attention}

\textbf{Motivation and mechanism.} Infini-attention combines masked local attention with a
bounded associative memory inside each block \citep{munkhdalai2024infini}.  With positive
feature map $\phi$, a segment updates
\begin{equation}
M_s=M_{s-1}+\phi(K_s)^\top V_s,\qquad
z_s=z_{s-1}+\sum_i\phi(k_{s,i}),
\end{equation}
and reads $\phi(Q_s)M_{s-1}/(\phi(Q_s)z_{s-1})$.  A learned gate mixes this compressive
read with dot-product attention over the current segment.  Memory size is independent of
total processed length, which makes streaming execution possible.

\begin{lstlisting}[caption={Associative-memory read and write}]
def infini(qf, kf, v, M, z, eps=1e-6):
    read = torch.einsum('btd,bdh->bth', qf, M)
    denom = torch.einsum('btd,bd->bt', qf, z).unsqueeze(-1)
    read = read / denom.clamp_min(eps)
    M = M + torch.einsum('btd,bth->bdh', kf, v)
    z = z + kf.sum(dim=1)
    return read, M.detach(), z.detach()
\end{lstlisting}
\methodflow{local segment}{kernel feature map}{bounded associative state}{gated local + memory read}{stream output}{memorygreen}

\textbf{Evidence and limitations.} The paper evaluates long-context language modeling,
one-million-token passkey retrieval, and 500K-token book summarization with 1B- and 8B-scale
models.  Metrics include perplexity, retrieval accuracy, summarization quality, and memory
scaling.  Bounded state is the strength and the limitation: many keys superpose, collisions
and normalization can obscure exact facts, and there is no explicit object version,
authorization boundary, or selective invalidation.

\textbf{PRA relation and direct comparison.} Infini-attention offers constant-size lossy
history; PRA offers selective identity with potentially unbounded backing storage and
bounded resident working set.  Match resident bytes and local window, then sweep number of
objects, repeated keys, exact-value queries, semantic summaries, and stale replacements.
PRA must include resolution/encoding latency.  A hybrid maintains cheap Infini state for
the stream while a collision, low-confidence query, or explicit handle faults in a named
PRA object.

\subsection{Titans and Test-Time Neural Memory}

\textbf{Motivation and mechanism.} Fixed-state recurrent models update activations, but
their update rule is usually shallow. Titans instead treats a neural network
$\mathcal M_{\theta_t}$ as long-term memory and updates its parameters while processing the
test sequence \citep{behrouz2025titans}. For a key--value association $(k_t,v_t)$, an
associative loss measures how poorly the current memory predicts the value,
\begin{equation}
\ell_t(\theta_t)=\|\mathcal M_{\theta_t}(k_t)-v_t\|_2^2,
\qquad s_t=\|\nabla_{\theta_t}\ell_t\|,
\end{equation}
and the gradient magnitude $s_t$ acts as a surprise signal. Momentum and forgetting terms
control how the parameter update incorporates surprising events. Titans combines this
neural long-term memory with precise short-window attention and data-independent persistent
memory. Its Memory-as-Context, Memory-as-Gate, and Memory-as-Layer variants differ in where
the long-term read enters the sequence model.

\begin{lstlisting}[caption={Conceptual surprise-weighted test-time memory update}]
def neural_memory_step(memory, key, value, optimizer_state):
    prediction = memory(key)
    loss = F.mse_loss(prediction, value)
    gradients = torch.autograd.grad(loss, memory.parameters())
    surprise = global_norm(gradients)
    optimizer_state = momentum_and_forgetting(optimizer_state, gradients, surprise)
    apply_test_time_update(memory, optimizer_state)
    return prediction, surprise
\end{lstlisting}
\methodflow{current token association}{neural memory loss}{surprise + momentum}{test-time parameter update}{long-term memory read}{memorygreen}

\textbf{Evidence and limitations.} Titans reports language modeling, common-sense
reasoning, genomics, time-series, and long-context retrieval experiments, including context
lengths beyond two million tokens. Ablations examine memory depth, momentum, forgetting,
convolution, persistent memory, and the three integration variants. These results motivate
test-time neural memory, but online parameter updates introduce sequential optimization,
state-reset, reproducibility, privacy, and contamination questions. The memory is content
addressed rather than an authoritative source store; forgetting or correcting one object
is not a simple URI invalidation.

\textbf{PRA relation and experiment.} Titans learns a compressed latent history online;
PRA resolves versioned external objects and transports selected evidence. Compare them on
four operations: repeated semantic patterns, exact quotations, correction of stale facts,
and access to an object never seen in the current stream. Cross both methods with state
reset, corrupted evidence, and repeated-query conditions. A hybrid could use Titans for
experience-dependent routing priors while PRA retains inspectable source identity. Such a
hybrid must report whether test-time updates improve object selection or merely adapt the
decoder independently of supplied content.

\section{Category III: Differentiable and Episodic External Memory}

These methods move state outside the ordinary recurrent or token representation while
keeping access near neural computation.  They introduce two questions that recur in PRA:
who writes the memory, and whether an address denotes a stable object or only a similar
vector.

\subsection{Memory Networks}

\textbf{Motivation and mechanism.} Memory Networks separate input encoding, memory
storage, addressing, and response production \citep{weston2015memory}.  Facts become slots
$m_i$; a query $q$ produces a content read
\begin{equation}
p_i=\operatorname{softmax}_i(q^\top m_i),\qquad
o=\sum_i p_i c_i,\qquad q' = q+o.
\end{equation}
Repeating the read implements multi-hop reasoning.  End-to-end variants learn embeddings
from question-answer supervision; earlier variants used stronger supporting-fact labels.

\begin{lstlisting}[caption={One differentiable memory hop}]
def memory_hop(query, address_slots, content_slots):
    weights = torch.softmax(query @ address_slots.transpose(-1, -2), dim=-1)
    read = weights @ content_slots
    return query + read, weights
\end{lstlisting}
\methodflow{query + stored facts}{content scores}{soft slot read}{repeat hops}{answer head}{memorygreen}

\textbf{Evidence, lineage, and failure modes.} The paper trail runs from supervised Memory
Networks through end-to-end memory networks and key-value memory networks.  Evaluations
include synthetic bAbI reasoning and question answering, using answer accuracy and
supporting-fact retrieval.  Multiple hops make compositional access inspectable, but soft
attention across a large flat store scales poorly; embeddings can exploit lexical shortcuts;
and overwrite, provenance, and slot lifecycle are under-specified for production knowledge.

\textbf{PRA relation and experiment.} Memory Networks establish the learned-read primitive
that PRA reuses, but PRA delegates object writes and identity to a resolver/cache.  Compare
flat learned slots with URI candidate restriction followed by the same within-candidate
attention.  Hold stored facts and read hops fixed; sweep memory size, paraphrase, distractors,
and wrong-object substitutions.  Report supporting-fact recall and counterfactual answer
change, not attention maps alone.  The expected trade-off is open semantic discovery versus
cheap exact restriction when a handle is available.

\subsection{Differentiable Neural Computer (DNC)}

\textbf{Motivation and mechanism.} The DNC couples a neural controller to a read/write
matrix with content lookup, temporal links, usage-based allocation, and differentiable
erase/add operations \citep{graves2016dnc}.  A write updates row $i$ as
\begin{equation}
M_t(i)=M_{t-1}(i)\odot(1-w_t^w(i)e_t)+w_t^w(i)a_t,
\end{equation}
while read weights mix content addressing with forward/backward traversal of a learned
temporal linkage.  It therefore supports both associative recall and ordered walks.

\begin{lstlisting}[caption={Core DNC erase/add write; allocation and link updates omitted}]
def dnc_write(memory, write_weight, erase, add):
    w = write_weight.unsqueeze(-1)
    return memory * (1 - w * erase.unsqueeze(1)) + w * add.unsqueeze(1)
\end{lstlisting}
\methodflow{controller state}{allocate/content address}{erase + write matrix}{content/link read heads}{controller output}{memorygreen}

\textbf{Evidence and limitations.} DNC experiments emphasize synthetic graph traversal,
question answering, and algorithmic tasks, comparing recurrent controllers and earlier
memory systems with exact task accuracy or sequence error. The machine can learn compact
algorithm-like behavior, yet many interacting gates make training brittle.  Soft access
is expensive at large slot counts, learned addresses lack stable external meaning, and
catastrophic write errors can corrupt later computation.

\textbf{PRA relation and experiment.} DNC learns both memory management and access;
PRA externalizes durable writes, versions, and permissions, narrowing learning to selection
and fusion.  Compare DNC, read-only DNC, and PRA on a graph task whose nodes have stable
IDs plus a variant with IDs hidden.  Match active memory bytes and controller parameters.
Test node substitution, temporal traversal, stale updates, and multi-tenant isolation.
A hybrid can use a small DNC as transient scratch space while PRA supplies immutable,
auditable backing objects.

\subsection{Memformer}

\textbf{Motivation and mechanism.} Memformer carries a fixed number of dynamic memory
slots across segments and introduces memory replay backpropagation so training need not
retain every old activation \citep{wu2020memformer}.  A segment attends to slots and local
tokens, then a write operation updates the slots:
\begin{equation}
S_{t+1}=\operatorname{Write}(S_t,H_t),\qquad |S_t|=N_s \text{ constant}.
\end{equation}
Replay reconstructs selected memory states during backward computation, trading extra
compute for lower activation storage.

\begin{lstlisting}[caption={Simplified fixed-slot segment update}]
def memformer_step(block, tokens, slots):
    joint = torch.cat([slots, tokens], dim=1)
    hidden = block(joint)
    new_slots = hidden[:, :slots.size(1)]
    token_hidden = hidden[:, slots.size(1):]
    return token_hidden, new_slots
\end{lstlisting}
\methodflow{fixed slots + segment}{joint attention}{update slots}{replay for gradients}{next segment}{memorygreen}

\textbf{Evidence, ablations, and failures.} The paper evaluates sequence modeling under
long inputs, reporting task quality, training memory, and speed against Transformer and
memory-augmented baselines.  Slot count, recurrence, and replay are key ablations.  Constant
space is attractive, but slots are lossy summaries without stable semantic names; useful
facts may be overwritten; replay adds complexity; and a slot cannot be independently
revoked or refreshed when its source changes.

\textbf{PRA relation and experiment.} Memformer provides a learned working set; PRA adds a
named backing store and explicit faults.  Use identical segment encoders and match resident
slots/KV bytes.  Ask questions whose evidence is either predictable during ingestion or
selected only after ingestion, and sweep time-to-query.  Compare slot reset/substitution
with PRA disabled/shuffled controls.  A tiered design stores Memformer summaries as cheap
URI metadata and faults in full layer-specific memory when summary confidence is low.

\subsection{Memorizing Transformer}

\textbf{Motivation and mechanism.} Memorizing Transformer records past key/value pairs in
a large non-differentiable episodic database and retrieves nearest keys for current queries
\citep{wu2022memorizing}.  For query $q$,
\begin{equation}
\mathcal I=\operatorname{kNN}(q,\{k_i\}),\qquad
o_{\rm mem}=\operatorname{softmax}(qK_{\mathcal I}^\top)V_{\mathcal I},
\end{equation}
which is gated with ordinary attention.  The store can span many batches and does not need
gradients through nearest-neighbor selection.

\begin{lstlisting}[caption={Episodic KV retrieval; an ANN index replaces the dense distance}]
def memorizing_read(q, key_db, value_db, k):
    dist = torch.cdist(q.float(), key_db.float())
    idx = dist.topk(k, largest=False).indices
    keys, values = key_db[idx], value_db[idx]
    score = torch.einsum('...d,...kd->...k', q, keys)
    return torch.einsum('...k,...kd->...d', score.softmax(-1), values)
\end{lstlisting}
\methodflow{current hidden query}{ANN over episodic keys}{top-$k$ past KV}{gated memory attention}{next-token state}{memorygreen}

\textbf{Evidence and failure modes.} Experiments span language-model datasets and model
scales up to billions of parameters, measuring perplexity versus memory size and retrieval
configuration.  The method demonstrates that latent episodic retrieval can improve a
Transformer without differentiating through a huge store.  Its risks are near-duplicate
memorization, index cost, stale states after model updates, cross-example contamination,
and weak source provenance.  Similarity can retrieve the right surface continuation while
failing compositional evidence use.

\textbf{PRA relation and experiment.} This is a close latent-memory precursor.  The main
difference is open kNN over anonymous stream states versus two-stage object identity and
within-object attention.  Compare (a) global episodic kNN, (b) exact URI restriction,
(c) semantic object routing then token kNN, and (d) oracle object restriction.  Hold total
indexed KVs and returned keys fixed.  Include near-duplicate continuation, multi-object
composition, wrong URI, and stale-model cache conditions, and report source-level recall
alongside perplexity.

\section{Category IV: Retrieval-Augmented Pretraining and Generation}

Retrieval augmentation places a corpus outside model parameters.  The decisive axes are
whether selection is supervised or latent, whether retrieved material enters as text or
hidden state, and whether retrieval occurs once, per chunk, or per generated token.

\subsection{kNN-LM}

\textbf{Motivation and mechanism.} kNN-LM stores training-context representations paired
with observed next tokens.  At inference it builds a distance-weighted neighbor distribution
and interpolates it with the parametric language model \citep{khandelwal2020knnlm}:
\begin{align}
p_{\rm kNN}(y\mid q)&\propto\sum_{i:y_i=y}
\exp(-d(q,k_i)/\tau),\\
p(y\mid q)&=\lambda p_{\rm kNN}(y\mid q)+(1-\lambda)p_{\rm LM}(y\mid q).
\end{align}
No retrieved passage is read compositionally; memory changes the final token distribution.

\begin{lstlisting}[caption={kNN-LM probability interpolation}]
def knn_lm(log_p_lm, q, keys, next_tokens, k, vocab, lam, tau):
    d, idx = torch.cdist(q, keys).topk(k, largest=False)
    weights = torch.softmax(-d / tau, dim=-1)
    p_knn = torch.zeros(*q.shape[:-1], vocab, device=q.device)
    p_knn.scatter_add_(-1, next_tokens[idx], weights)
    return torch.logaddexp(log_p_lm + math.log1p(-lam),
                           p_knn.clamp_min(1e-12).log() + math.log(lam))
\end{lstlisting}
\methodflow{LM context vector}{kNN datastore}{neighbor token labels}{distribution interpolation}{next token}{retrievalblue}

\textbf{Evidence, lineage, and limitations.} The original work reports perplexity on
WikiText-103 and domain-adaptation settings, comparing the base LM and cache/retrieval
baselines.  Improvements without parameter updates motivate later datastore adaptation,
pruning, and efficient kNN work.  Storage and ANN latency can be large; interpolation needs
calibration; privacy risks follow from training-example memorization; and the mechanism is
strongest for continuation similarity rather than multi-fact reasoning.

\textbf{PRA relation and comparison.} kNN-LM fuses at the output distribution; PRA fuses
rich memory within hidden computation.  Compare them with the same corpus encoder and
retrieval budget on near-duplicate prediction versus questions requiring composition.
Report perplexity, calibration, exact-match, datastore bytes, and latency.  Add shuffled
token labels for kNN-LM and shuffled objects for PRA.  If PRA only matches kNN-LM on
near-copy tasks at higher cost, its richer transport is not justified.

\subsection{REALM}

\textbf{Motivation and mechanism.} REALM jointly pretrains a dense retriever and a masked
language model over an external corpus \citep{guu2020realm}.  With latent document $z$,
\begin{equation}
p(y\mid x)=\sum_{z\in\operatorname{TopK}(x)}p_\eta(z\mid x)
p_\theta(y\mid x,z),
\end{equation}
where approximate maximum inner-product search supplies candidates and the language loss
provides a learning signal to the retriever.  Retrieved text is encoded jointly with the
masked input, so the knowledge store can change without retraining all model parameters.

\begin{lstlisting}[caption={REALM-style latent-document marginalization}]
scores, doc_ids = retriever.topk(query_encoder(masked_input), k)
log_p_doc = scores.log_softmax(-1)
log_p_y = torch.stack([reader(masked_input, corpus[i]) for i in doc_ids], -1)
loss = -torch.logsumexp(log_p_doc + log_p_y, dim=-1).mean()
\end{lstlisting}
\methodflow{masked text query}{dense corpus retrieval}{encode top documents}{latent marginalization}{token/QA prediction}{retrievalblue}

\textbf{Evidence and failure modes.} REALM evaluates open-domain QA and knowledge-intensive
pretraining, comparing retrieval-free language models and retrieval pipelines with exact
match and retrieval diagnostics.  Its important legacy is end-to-end retriever learning
from the task objective.  Index refresh is expensive, top-$k$ truncation biases the latent
sum, corpus snapshots become stale, and retrieval may learn shortcuts when answers are
predictable locally.

\textbf{PRA relation and experiment.} PRA should borrow REALM's trainable selection rather
than rely on fixed summary cosine scores.  REALM materializes retrieved text in an encoder;
PRA can reuse per-object layer KVs and explicit handles.  Cross oracle/learned retrieval
with text/cross-attention transport on the same corpus.  Match retrieved tokens and index
queries; evaluate locally sufficient and reference-required subsets, shuffled documents,
random handles, warm caches, and index refresh after source updates.

\subsection{Retrieval-Augmented Generation (RAG)}

\textbf{Motivation and mechanism.} RAG combines a dense retriever with a seq2seq generator.
RAG-Sequence holds a document mixture for a whole output; RAG-Token permits the latent
document to vary by token \citep{lewis2020rag}:
\begin{equation}
p(y\mid x)=\prod_t\sum_{z\in\operatorname{TopK}(x)}
p_\eta(z\mid x)p_\theta(y_t\mid x,z,y_{<t}).
\end{equation}
In common systems practice, retrieved passages are concatenated into a prompt rather than
exactly marginalized, but the selection/materialization distinction is the same.

\begin{lstlisting}[caption={Minimal retrieve-then-read RAG pipeline}]
def rag_answer(question, retriever, generator, k):
    hits = retriever.search(question, k=k)
    context = "\n\n".join(f"[{h.id}] {h.text}" for h in hits)
    return generator(f"Sources:\n{context}\nQuestion: {question}\nAnswer:")
\end{lstlisting}
\methodflow{question}{dense/sparse top-$k$}{passage text}{reader or prompt fusion}{answer + citations}{retrievalblue}

\textbf{Evidence, descendants, and failures.} RAG evaluates Natural Questions, TriviaQA,
WebQuestions, CuratedTREC, MS MARCO-style knowledge tasks, and generation, comparing DPR,
closed-book, and retrieval-augmented baselines with exact match and generation/diversity
metrics.  It seeded extensive work on reranking, query rewriting, corrective and agentic
RAG.  Failure modes include chunk boundary errors, retrieval miss, distractor sensitivity,
prompt-token cost, citation mismatch, and stale indexes.  A fluent answer does not prove
that retrieved evidence caused the output.

\textbf{PRA relation and direct comparison.} Both keep knowledge external; they differ in
timing and materialization.  RAG usually chooses text before generation, while PRA can
expose named latent objects by layer or step.  Use one corpus, retriever, generator backbone,
and evidence budget.  Compare prompt text, encoder cross-attention, cached object KVs, and
iterative retrieval.  Report task quality, citation precision/recall, valid/disabled/shuffled
effects, cold/warm latency, tokenization/encoding cost, and behavior as queries over the same
documents repeat.

\subsection{Atlas}

\textbf{Motivation and mechanism.} Atlas asks whether retrieval-augmented pretraining can
make knowledge-intensive tasks data efficient, rather than treating retrieval only as an
inference add-on \citep{izacard2022atlas}. A dual-encoder retriever selects passages and a
Fusion-in-Decoder reader encodes each question--passage pair separately before the decoder
attends across all encoded passages. Retriever supervision can come from reader attention
or leave-one-out evidence estimates. Abstractly,
\begin{equation}
z_{1:k}=\operatorname{TopK}_{z\in\mathcal D}s_\eta(q,z),\qquad
p(y\mid q,z_{1:k})=\operatorname{FiD}_\theta(q,z_{1:k}).
\end{equation}
The external index can be updated without storing all knowledge in model parameters, while
joint pretraining teaches the retriever and reader a shared task interface.

\begin{lstlisting}[caption={Atlas-style retrieve, encode separately, and fuse in the decoder}]
def atlas_forward(question, retriever, reader, corpus, k):
    passages = corpus.fetch(retriever.topk(question, k))
    encoded = [reader.encode(question, passage) for passage in passages]
    return reader.decode(torch.cat(encoded, dim=1))
\end{lstlisting}
\methodflow{few-shot question}{jointly pretrained retriever}{separate passage encodings}{decoder fusion}{knowledge-task output}{retrievalblue}

\textbf{Evidence and limitations.} Atlas evaluates few-shot and full-data learning on
Natural Questions, KILT tasks, MMLU, and related knowledge-intensive benchmarks. The paper
also varies index content and retrieval-training objectives. Its contribution is strongest
at the training interface: retrieval provides updateable evidence and sample-efficient task
adaptation. It still pays to retrieve and encode a passage set before generation, and
reader-derived retriever targets can reinforce the reader's own shortcuts.

\textbf{PRA relation and experiment.} Atlas supplies a stronger pretrained retrieval
baseline than prompt-only RAG. A direct comparison should initialize both systems from the
same corpus and language backbone, then vary the number of task examples. Atlas materializes
passage encodings; PRA materializes selected per-layer chunk K/V. Match retrieved tokens,
pretraining data, and online encoding cost. Oracle passages test transport, while
valid-versus-shuffled evidence and index-refresh tests measure content and updateability.

\subsection{Self-RAG}

\textbf{Motivation and mechanism.} Fixed-$k$ RAG retrieves even when parametric knowledge
is sufficient and may accept irrelevant evidence without an explicit check. Self-RAG trains
the generator to emit reflection tokens that control retrieval and assess passage relevance,
support, and response quality \citep{asai2024selfrag}. A simplified policy alternates
\begin{equation}
a_t\in\{\textsc{retrieve},\textsc{continue}\},\qquad
r_t=\operatorname{critique}(q,z_t,y_{<t}),
\end{equation}
then uses critique scores to filter or rank candidate continuations. Retrieval triggering
and evidence criticism therefore become learned generation decisions rather than fixed
pipeline settings.

\begin{lstlisting}[caption={Simplified adaptive retrieval and reflection loop}]
while not finished:
    action = model.next_control_token(question, draft)
    if action == RETRIEVE:
        passages = retriever.search(question + draft)
        passages = rank_by_reflection(model, question, draft, passages)
    draft += model.generate_segment(question, draft, passages)
return draft
\end{lstlisting}
\methodflow{question + partial answer}{retrieve decision}{passage relevance critique}{supported segment generation}{response critique}{retrievalblue}

\textbf{Evidence and limitations.} Self-RAG evaluates open-domain QA, reasoning, fact
verification, long-form factuality, and citation behavior with 7B- and 13B-scale models.
Reflection tokens make retrieval frequency controllable, but their labels depend on a
teacher and critic pipeline. A plausible self-critique can share the generator's errors,
and discrete retrieve/generate cycles still incur text and serial decoding cost.

\textbf{PRA relation and experiment.} Self-RAG is the closest retrieval analogue to PRA's
proposed progressive trigger. Compare a reflection-token trigger, a latent PRA gate, and a
shared oracle trigger while holding the retriever and evidence fixed. Measure trigger
precision/recall, unnecessary accesses, evidence support, serial latency, and content
counterfactuals. A hybrid can expose PRA route confidence and provenance to a Self-RAG
critic, but the critic's approval must not replace removal and substitution tests.

\subsection{RETRO}

\textbf{Motivation and mechanism.} RETRO retrieves neighbors for fixed-size chunks from a
massive token database, encodes each neighbor and its continuation, and interleaves chunked
cross-attention in the decoder \citep{borgeaud2022retro}.  For input chunk $x_c$ and
neighbors $N_c$,
\begin{equation}
N_c=\operatorname{kNN}(E(x_c),\mathcal D),\qquad
H_c'=H_c+\operatorname{CrossAttn}(H_c,\operatorname{Enc}(N_c)).
\end{equation}
Retrieval happens at chunk boundaries, bringing non-parametric data closer to layer
computation than prompt-only RAG.

\begin{lstlisting}[caption={RETRO-style chunk retrieval and cross-attention}]
for chunk in split(tokens, chunk_len):
    ids = index.search(chunk_encoder(chunk), k=neighbors)
    neighbor_states = encoder(corpus.with_continuations(ids))
    h = decoder_chunk(chunk, retrieved=neighbor_states)  # selected layers cross-attend
\end{lstlisting}
\methodflow{input chunk}{nearest corpus chunks}{encode neighbor + continuation}{chunked cross-attention}{decoder layers}{retrievalblue}

\textbf{Evidence, ablations, and limitations.} RETRO trains models with a database on the
order of trillions of tokens and compares perplexity and downstream behavior against larger
parametric Transformers.  Retrieval neighbors, database scale, cross-attention placement,
and corpus overlap are central ablations.  Results support non-parametric scaling, but the
system inherits corpus/index cost, possible benchmark leakage, chunk-boundary mismatch,
and retrieval latency; source identity is incidental rather than an explicit user-facing
contract.

\textbf{PRA relation and experiment.} RETRO is one of PRA's closest architectural ancestors:
both transport external representations through cross-attention.  PRA adds explicit handles,
object boundaries, layer-specific reusable cache identity, progressive admission, and causal
reference controls.  Under one corpus and backbone, compare no handle, noisy handle, exact
handle, and global semantic retrieval; keep returned tokens and PRA layers fixed.  Measure
selection recall, quality, per-stage latency, cache reuse, source corruption, and multi-object
composition.  Oracle selection isolates whether PRA's transport works before router claims.

\section{Category V: Attention-Level Retrieval and KV Runtime Management}

This category is easy to conflate because every method manipulates latent states. The
systems nevertheless perform different operations. Unlimiformer and Landmark Attention
select semantic states or blocks. PagedAttention virtualizes physical allocation.
StreamingLLM, H$_2$O, Scissorhands, SnapKV, FastGen, PyramidKV, and Quest reduce which
already-produced states are loaded or retained. MiniCache compresses across layers, while
CacheBlend repairs reusable chunk caches that were encoded outside their new prompt
position. Selection, allocation, retention, representation compression, and reuse should
therefore be measured separately.

\subsection{Unlimiformer}

\textbf{Motivation and mechanism.} Unlimiformer lets a pretrained encoder-decoder process
arbitrarily segmented input, indexes all encoder hidden states, and replaces dense decoder
cross-attention with exact attention over MIPS-selected states \citep{bertsch2023unlimiformer}.
Because the cross-attention score is an inner product, each decoder query can retrieve its
largest candidate logits without first forming the full query--key matrix:
\begin{equation}
\mathcal I(q)=\operatorname{TopK}_i(q^\top k_i),\qquad
o=\operatorname{softmax}(qK_{\mathcal I}^\top)V_{\mathcal I}.
\end{equation}

\begin{lstlisting}[caption={Token-state MIPS followed by exact attention}]
def unlimiformer_cross_attention(q, index, value_store, k):
    scores, ids = index.search(q, k)
    values = value_store[ids]
    return torch.einsum('...k,...kd->...d', scores.softmax(-1), values)
\end{lstlisting}
\methodflow{segmented long input}{index encoder states}{per-query top-$k$ MIPS}{exact sparse cross-attention}{decoder output}{retrievalblue}

\textbf{Evidence and limitations.} The paper evaluates long-input encoder-decoder tasks,
especially summarization, against truncated, segmented, and long-context baselines using
ROUGE and memory/length measures.  It also studies training-free wrapping and fine-tuning.
Selection occurs at token granularity and can miss jointly useful low-ranked tokens; the
index is tied to the current encoded input; ANN setup and query cost remain; and source
provenance must be reconstructed from token offsets.

\textbf{PRA relation and comparison.} Unlimiformer and PRA meet at sparse latent attention.
Unlimiformer indexes all token states of the current long input, while PRA first selects
named objects and may persist their per-layer KVs.  Compare full cross-attention,
Unlimiformer token MIPS, PRA object routing plus dense within-object attention, and a
hierarchical object-to-token cascade.  Match returned KVs, encoder FLOPs, and peak bytes;
sweep object size, distractors, repeated queries, and object-level corruption.  The cascade
tests whether explicit identity improves efficiency without sacrificing token precision.

\subsection{Landmark Attention}

\textbf{Motivation and mechanism.} Ordinary sparse attention requires a fixed graph or an
external retriever. Landmark Attention instead trains the model's own attention to use one
special landmark token as the representative key for each token block
\citep{mohtashami2023landmark}. At inference, a query first scores landmarks, retrieves the
top blocks, and then attends to token K/V inside those blocks. If $\lambda_b$ is the
landmark for block $b$,
\begin{equation}
\mathcal B(q)=\operatorname{TopK}_b(q^\top k_{\lambda_b}),\qquad
o=\operatorname{Attn}(q,K_{\mathcal B(q)},V_{\mathcal B(q)}).
\end{equation}
The method preserves random access to distant blocks while avoiding attention over every
token. Its optional context-miss token also suggests a learned decision not to retrieve.

\begin{lstlisting}[caption={Landmark routing followed by block attention}]
def landmark_attention(query, landmark_keys, block_keys, block_values, k):
    block_ids = (query @ landmark_keys.T).topk(k).indices
    keys = torch.cat([block_keys[i] for i in block_ids], dim=-2)
    values = torch.cat([block_values[i] for i in block_ids], dim=-2)
    return scaled_dot_product_attention(query, keys, values)
\end{lstlisting}
\methodflow{token blocks + landmarks}{query scores landmarks}{top-$k$ block fetch}{token-level block attention}{random-access output}{retrievalblue}

\textbf{Evidence and limitations.} The paper evaluates PG19 language modeling and passkey
retrieval, including a LLaMA-7B adaptation beyond its original context. It studies block
retrieval granularity and a context-miss threshold. Landmark keys are learned in the same
attention space as their consumers, but every block must be ingested and assigned a
landmark. Blocks remain positional rather than stable versioned objects, and a wrong
landmark decision hides every token inside the missed block.

\textbf{PRA relation and experiment.} Landmark Attention is the closest token-stream
analogue to PRA's corrected chunk gists. Both rank a small projected-key representative and
transport detailed token K/V after selection. The differences are identity, lifecycle,
and scope: landmarks index blocks of an ingested stream, whereas PRA gists belong to
resolved URI objects that may be cached and invalidated. Compare both with identical chunk
boundaries, top-$k$, and attention projections; then vary object reuse, reorder blocks,
change versions, and measure cold versus warm construction. This experiment can reveal
whether explicit object identity adds value beyond learned block routing.

\subsection{PagedAttention}

\textbf{Motivation and mechanism.} PagedAttention is a serving-memory design in vLLM, not
a semantic attention rule.  It partitions a sequence's KV cache into fixed-size physical
blocks addressed through a logical block table, reducing fragmentation and enabling
copy-on-write sharing \citep{kwon2023vllm}.  For logical token position $t$,
\begin{equation}
b=\lfloor t/B\rfloor,\quad o=t\bmod B,\quad
(K_t,V_t)=\operatorname{blockTable}[b][o].
\end{equation}

\begin{lstlisting}[caption={Logical-to-physical KV lookup}]
def paged_lookup(logical_positions, block_table, blocks, block_size):
    logical_block = logical_positions // block_size
    offset = logical_positions % block_size
    physical_block = block_table[logical_block]
    return blocks[physical_block, offset]
\end{lstlisting}
\methodflow{logical sequence KV}{block table}{non-contiguous physical blocks}{paged attention kernel}{batched decoding}{localgray}

\textbf{Evidence, systems lineage, and failures.} vLLM evaluates serving throughput,
latency, memory utilization, and scheduling across model sizes and workloads against
contemporary inference systems.  Block size, batching, sharing, and scheduling are systems
variables rather than accuracy ablations.  Paging cannot decide which knowledge is useful;
small blocks increase metadata, large blocks reintroduce internal fragmentation, and shared
physical storage must preserve isolation and copy-on-write semantics.

\textbf{PRA relation and experiment.} The designs compose: PRA supplies semantic object
identity and PagedAttention supplies allocation.  Map an immutable $(u,v,\ell,m,p)$ object
to reference-counted KV blocks and maintain per-sequence block tables/masks.  Compare
contiguous and paged PRA caches under cold, warm, and authorized cross-request reuse.
Measure bytes allocated, fragmentation, blocks shared, avoided encoding, time to first token,
throughput, eviction, and a confidentiality test where another sequence's private cache
must leave the target logits invariant.

\subsection{StreamingLLM}

\textbf{Motivation and mechanism.} StreamingLLM observes that autoregressive models often
assign disproportionate attention to initial ``sink'' tokens even when those tokens are
not semantically important.  Retaining a few initial KVs plus a rolling recent window
allows stable streaming beyond the training context \citep{xiao2023streamingllm}:
\begin{equation}
\mathcal K_t=\{0,\ldots,s-1\}\cup\{t-w+1,\ldots,t\}.
\end{equation}
Middle-history KVs are evicted; no semantic retrieval or reconstruction is attempted.

\begin{lstlisting}[caption={Attention-sink plus rolling-window retention}]
def streaming_keep(length, sinks=4, window=1024):
    first = torch.arange(min(sinks, length))
    recent = torch.arange(max(sinks, length-window), length)
    return torch.unique(torch.cat([first, recent]))
\end{lstlisting}
\methodflow{growing token stream}{keep sink tokens}{keep recent window}{evict middle KV}{stable streaming decode}{localgray}

\textbf{Evidence and failure modes.} The paper tests long streaming language modeling and
generation with LLaMA-style models, comparing sliding windows and cache baselines through
perplexity and long-run behavior.  Sink tokens repair an attention-distribution pathology;
they do not preserve arbitrary facts.  Middle-history details are irrecoverable, sink
behavior depends on model training, and the method is not a solution for non-chronological
external knowledge.

\textbf{PRA relation and experiment.} StreamingLLM manages the local chronological cache;
PRA retrieves named backing objects.  Combine them by keeping sinks/recent discourse locally
and dereferencing remote evidence.  Under a fixed total KV budget, compare rolling window,
sink+window, PRA, and the hybrid.  Sweep query distance, explicit handles, and whether lost
facts are still present in the backing store.  Charge PRA faults and test shuffled content;
test StreamingLLM with sink removal and exact middle-history probes.

\subsection{\texorpdfstring{H$_2$O}{H2O}: Heavy-Hitter Oracle}

\textbf{Motivation and mechanism.} H$_2$O treats KV eviction as a dynamic submodular
selection problem.  It retains recent tokens and ``heavy hitters'' with large accumulated
attention contribution \citep{zhang2023h2o}.  A simplified score is
\begin{equation}
I_i^{(t)}=I_i^{(t-1)}+\sum_h A_{h,t,i},\qquad
\mathcal K_t=\operatorname{Recent}_r\cup\operatorname{TopK}_h(I^{(t)}).
\end{equation}
The policy reacts to observed decoder attention rather than using a static window.

\begin{lstlisting}[caption={Educational heavy-hitter update}]
def h2o_keep(importance, attn_to_past, recent_ids, heavy_budget):
    importance[:attn_to_past.numel()] += attn_to_past.detach()
    heavy = importance.topk(heavy_budget).indices
    return torch.unique(torch.cat([heavy, recent_ids])), importance
\end{lstlisting}
\methodflow{past stream KV}{accumulate attention importance}{heavy + recent selection}{evict low-score KV}{memory-bounded decode}{localgray}

\textbf{Evidence and limitations.} H$_2$O evaluates OPT, LLaMA, and GPT-NeoX across language
and generation tasks, measuring quality preservation, cache size, throughput, and latency
against inference systems and eviction baselines.  Its reported systems gains are conditional
on model, hardware, batch, cache fraction, and implementation.  Accumulated attention is a
lagging proxy, early errors are irreversible, low-attention setup facts can matter later,
and raw attention mass is not causal importance.

\textbf{PRA relation and experiment.} H$_2$O decides which states from one active stream
remain; PRA decides which external objects enter.  Use H$_2$O within each selected PRA object
or across the resident object working set, but preserve URI/version metadata for every KV.
Compare full object KVs, H$_2$O-compressed objects, object-level eviction, and both tiers.
Sweep delayed-use facts, distractors, object count, and warm-cache reuse.  Removal/substitution
tests should verify that retained heavy hitters, rather than generic branch activation,
cause correct answers.

\subsection{Scissorhands}

\textbf{Motivation and mechanism.} Scissorhands builds on the ``persistence of importance''
hypothesis: tokens important in earlier attention tend to remain important later.  It
maintains importance statistics and discards KVs predicted to stay unimportant
\citep{liu2023scissorhands}.  At an eviction step,
\begin{equation}
\mathcal K=\operatorname{Recent}_r\cup
\operatorname{TopK}_{B-r}(\operatorname{score}_{\rm history}(i)).
\end{equation}
Policies can be applied per head/layer and periodically rather than at every token.

\begin{lstlisting}[caption={Periodic history-score pruning}]
def scissorhands(kv, score, recent, budget):
    protected = torch.arange(max(0, len(score)-recent), len(score))
    old = torch.arange(0, max(0, len(score)-recent))
    keep_old = old[score[old].topk(max(0, budget-len(protected))).indices]
    keep = torch.cat([keep_old, protected])
    return kv.index_select(-2, keep), score.index_select(0, keep)
\end{lstlisting}
\methodflow{full past KV}{historical importance}{protect recent tokens}{periodic token pruning}{smaller active cache}{localgray}

\textbf{Evidence and failure modes.} The work evaluates cache reduction across language
models/tasks using perplexity or task quality plus memory and throughput.  Comparisons with
window and heuristic eviction establish the value of history-aware scores.  Importance
persistence can fail under topic shifts or late references; dropped tokens cannot be
reconstructed; per-head policies complicate kernels; and attention history remains a
correlational diagnostic.

\textbf{PRA relation and experiment.} Scissorhands compresses already materialized stream
history, whereas PRA can reload a stable object after eviction.  Compare irreversible token
eviction, eviction plus raw backing-text reload, and eviction plus cached PRA object KVs.
Match total bytes and include late queries after long distractor gaps.  Report reload rate,
quality, end-to-end latency, and counterfactual content effects.  Object boundaries may
provide a cleaner eviction unit than anonymous tokens, but that is a testable hypothesis.

\subsection{SnapKV}

\textbf{Motivation and mechanism.} SnapKV compresses a long prompt's KV cache before
generation by using attention from a recent observation window to identify prompt positions
likely to matter during decoding, then pooling nearby selections \citep{li2024snapkv}.  A
schematic per-head rule is
\begin{equation}
s_i=\operatorname{pool}\left(\sum_{q\in W_{\rm obs}}A_{q,i}\right),\qquad
\mathcal K=\operatorname{TopK}_B(s).
\end{equation}
Unlike online eviction, it takes a ``snapshot'' of prompt relevance before the answer grows.

\begin{lstlisting}[caption={Observation-window prompt selection}]
def snapkv_scores(attn, observation, pool=7):
    # attn: [heads, query, key]; exclude generated/recent keys as appropriate
    score = attn[:, -observation:, :].sum(dim=1).unsqueeze(1)
    return F.max_pool1d(score, kernel_size=pool, stride=1, padding=pool//2).squeeze(1)
\end{lstlisting}
\methodflow{long prompt KV}{recent observation window}{per-head attention votes}{cluster + retain positions}{compressed generation KV}{localgray}

\textbf{Evidence and limitations.} SnapKV evaluates long-context QA/retrieval and generation
across instruction-tuned LLMs, measuring task quality, prompt-cache reduction, and speed
against eviction/compression baselines.  Its main assumption is that the observation window
predicts future importance.  Later reasoning can change the relevant region; prompt
positions lose object/version identity; selection errors are irreversible unless the source
remains elsewhere; and head-wise layouts complicate reuse.

\textbf{PRA relation and experiment.} SnapKV is query-aware token compression after eager
prompt ingestion; PRA is object admission before or during latent use.  Compare full prompt,
SnapKV, PRA, and PRA with SnapKV-compressed object KVs.  Match resident bytes and include
questions whose relevance is evident initially versus revealed after intermediate reasoning.
Measure prefill, selection, cold encoding, decode latency, exact evidence retention, and
valid/shuffled object effects.  A multi-fidelity PRA cache could store a SnapKV snapshot and
escalate to full object KVs on low confidence.

\subsection{FastGen and PyramidKV: Head- and Layer-Adaptive Budgets}

\textbf{Motivation and mechanism.} A uniform cache budget assumes that every attention
head and layer uses history in the same way. FastGen profiles each head and assigns a
retention policy that matches its observed structure: local heads keep a window,
special-token heads keep designated positions, and broadly distributed heads retain more
of the cache \citep{ge2023fastgen}. PyramidKV operates along a different axis. It reports
that attention is broad in lower layers and more concentrated in higher layers, then uses
a decreasing layer budget \citep{cai2024pyramidkv}. A generic allocation is
\begin{equation}
B_{\ell,h}=B_{\rm total}\,
\frac{w_{\ell,h}}{\sum_{j,g}w_{j,g}},
\end{equation}
where the weight $w_{\ell,h}$ comes from head profiling, layer depth, or both. The equation
does not prescribe either paper's exact policy; it exposes their shared principle that
cache capacity should follow attention structure rather than tensor shape.

\begin{lstlisting}[caption={Composing head policy with a layer-dependent cache budget}]
for layer, heads in enumerate(model.attention_heads):
    layer_budget = pyramid_budget(layer, model.num_layers, total_budget)
    for head in heads:
        policy = profiled_policy(head)  # local, special-token, or broad
        head.cache = policy.retain(head.cache, budget=layer_budget[head.id])
\end{lstlisting}
\methodflow{full per-layer KV}{profile head behavior}{allocate pyramidal budgets}{policy-specific retention}{non-uniform cache}{localgray}

\textbf{Evidence and limitations.} FastGen evaluates training-free adaptive compression
across language-model tasks and reports memory savings with small quality changes under its
tested settings. PyramidKV evaluates LongBench and needle retrieval, including aggressive
retention fractions. Both depend on attention patterns transferring across inputs. A head
classified as local can become globally important on another task, and a depth trend does
not identify which individual evidence tokens must survive. Specialized per-head layouts
can also reduce kernel regularity.

\textbf{PRA relation and experiment.} These methods can compress each selected PRA object's
token K/V after object and chunk routing. Compare uniform and adaptive budgets under the
same total bytes, with per-layer routing recall and content counterfactuals. The critical
failure test delays a reference until a head or layer previously classified as low-value
must retrieve it. A successful composition would reduce resident bytes without changing
which URI and chunk caused the answer.

\subsection{Quest: Query-Aware Page Selection}

\textbf{Motivation and mechanism.} Retention policies decide what survives before the
future query is known. Quest retains paged K/V metadata and decides which pages to load for
the current query \citep{tang2024quest}. For each key dimension, a page stores extrema
$k^{\min}_p$ and $k^{\max}_p$. The query signs choose an upper bound on the page's possible
inner product,
\begin{equation}
U(q,p)=\sum_j \max(q_j k^{\min}_{p,j},q_j k^{\max}_{p,j}),
\qquad \mathcal P(q)=\operatorname{TopK}_p U(q,p).
\end{equation}
Only selected pages are transferred for exact attention. The method therefore reduces
memory bandwidth without first discarding the full backing cache.

\begin{lstlisting}[caption={Quest-style query bound over KV pages}]
def quest_page_scores(query, page_key_min, page_key_max):
    lower = query * page_key_min
    upper = query * page_key_max
    return torch.maximum(lower, upper).sum(dim=-1)

pages = quest_page_scores(q, key_min, key_max).topk(page_budget).indices
output = exact_attention(q, load_keys(pages), load_values(pages))
\end{lstlisting}
\methodflow{paged full backing KV}{query + page extrema}{upper-bound top-$k$}{load selected pages}{exact sparse attention}{retrievalblue}

\textbf{Evidence and limitations.} Quest reports self-attention speedups and end-to-end
latency improvements on long-context tasks while preserving quality in the evaluated
configurations. Its page bounds are inexpensive and query dependent, but loose bounds load
irrelevant pages, tight budgets can omit diffuse evidence, and page layout determines
selection granularity. The full cache still occupies a backing tier even when HBM traffic
falls.

\textbf{PRA relation and experiment.} Quest performs physical page routing inside an
already materialized context; PRA performs semantic object and chunk routing before
transport. The two-stage composition is natural: PRA selects URI chunks and Quest selects
pages inside their token K/V. Compare flat Quest over all candidate objects, hierarchical
PRA-to-Quest, and PRA with dense selected chunks. Match page metadata, loaded bytes, and
top-$k$ tokens. Report both semantic recall and actual bandwidth, because a good router that
loads loose pages may not improve the systems frontier.

\subsection{MiniCache and CacheBlend: Depth Compression and Reusable Chunks}

\textbf{Motivation and mechanism.} Token eviction compresses the sequence dimension, but
KV redundancy also appears across depth and repeated prompts. MiniCache observes similarity
between adjacent middle-to-late layers and merges corresponding states after separating
vector direction from magnitude \citep{liu2024minicache}. Distinct state pairs can remain
unmerged. CacheBlend addresses a different reuse problem: a document chunk encoded alone
does not contain the cross-chunk interactions it would receive in a new RAG prompt. It
reuses offline chunk K/V but selectively recomputes important tokens to repair the blended
context \citep{yao2024cacheblend}.

\begin{lstlisting}[caption={Two orthogonal cache transformations}]
# Depth compression: share similar adjacent-layer states.
merged = merge_direction_preserve_norm(kv[layer], kv[layer + 1], keep_distinct=True)

# Reuse repair: blend cached chunks, then recompute selected contextual tokens.
prompt_kv = concatenate(cached_chunk_kv)
repair_ids = select_recompute_tokens(prompt_kv, query)
prompt_kv[repair_ids] = model.prefill_tokens(prompt, repair_ids)
\end{lstlisting}
\methodflow{layer or chunk caches}{detect reusable structure}{merge or concatenate}{retain/recompute exceptions}{smaller prefill cost}{localgray}

\textbf{Evidence and limitations.} MiniCache evaluates several LLaMA, Mistral, Phi, and
Mixtral variants across language and long-context benchmarks, reporting compression,
throughput, and quality. CacheBlend evaluates RAG serving with several open models and
datasets, reporting time to first token and throughput against full recomputation and
prefix reuse. Cross-layer similarity is model dependent, while chunk reuse is position and
context dependent. In both cases, approximate equivalence must be tested on compositional
questions, not inferred from cache similarity alone.

\textbf{PRA relation and experiment.} MiniCache can reduce the depth cost of a PRA object;
CacheBlend is a direct baseline for PRA's reusable-object systems claim. Compare full
per-request encoding, prefix-only reuse, CacheBlend repair, position-independent PRA
cross-attention, and PRA with depth-compressed object K/V. Use reordered documents,
multi-document composition, stale versions, and shared-cache workloads. Measure answer
quality, repair fraction, encoding avoided, bytes transferred, time to first token, and
isolation. PRA should not claim reusable K/V advantage unless it beats these stronger
reuse baselines at matched quality.

\subsection{KV management design map}

The methods above can be compared by the unit they control and by whether omitted detail
still exists in a backing tier. That distinction predicts recoverability: page selection
can load another page on the next query, whereas destructive token eviction cannot recover
a dropped state unless raw content is re-encoded.

\begin{table}[htbp]
\centering\scriptsize
\caption{KV-management mechanisms control different logical and physical units.}
\begin{tabular}{>{\raggedright\arraybackslash}p{2.5cm}>{\raggedright\arraybackslash}p{2.0cm}>{\raggedright\arraybackslash}p{1.9cm}>{\raggedright\arraybackslash}p{2.4cm}>{\raggedright\arraybackslash}p{1.8cm}>{\raggedright\arraybackslash}p{2.8cm}}
\toprule
Method & Controlled unit & Decision time & Backing detail & Training change & Principal risk \\\midrule
PagedAttention & physical block & allocation & full K/V remains & no & fragmentation or isolation bug \\
StreamingLLM & token positions & decode & evicted middle absent & no & unrecoverable old detail \\
H$_2$O & token positions & decode & evicted tokens absent & no & lagging attention score \\
Scissorhands & token/head K/V & periodic decode & evicted tokens absent & no & importance changes later \\
SnapKV & prompt token/head & before generation & source may remain external & no & observation window predicts poorly \\
FastGen & head policy & profile/prefill & policy dependent & no fine-tuning & profile fails to transfer \\
PyramidKV & layer budget & prefill & omitted tokens absent & no & depth heuristic misses evidence \\
Quest & K/V page & each query & full backing K/V remains & no & loose bounds or page miss \\
MiniCache & adjacent layers & prefill/decode & selected exceptions remain & no & cross-layer mismatch \\
CacheBlend & reusable text chunk & prefill & raw text and offline K/V remain & no & insufficient contextual repair \\
Landmark Attention & token block & each query & block K/V remains & yes & wrong landmark hides block \\
PRA & URI chunk, then token K/V & layer/step & raw/versioned object may remain & yes/runtime policy & wrong route, stale identity \\
\bottomrule
\end{tabular}
\end{table}

The map exposes three promising compositions. PagedAttention can allocate any of the
logical caches. Quest can reduce bandwidth inside a selected Landmark or PRA block.
MiniCache or an adaptive budget can compress the resulting resident tensors. CacheBlend,
by contrast, is a direct alternative when the goal is reuse of text chunks inside ordinary
self-attention. End-to-end comparisons must include all metadata, repair, transfer, and
re-encoding costs; nominal cache bytes alone do not identify the faster system.

\section{Category VI: Structural and Hierarchical Indexing}

Flat chunk retrieval treats document structure as noise.  Hierarchical and graph approaches
make sentences, sections, entities, communities, or claims explicit intermediate units.
Their advantage is not merely fewer tokens: the index supplies a reasoning topology and a
natural provenance coordinate system.

\subsection{Hierarchical Attention Networks (HAN)}

\textbf{Motivation and mechanism.} HAN mirrors document composition: a bidirectional word
encoder and attention pool create sentence vectors, then a sentence encoder and attention
pool create a document vector \citep{yang2016han}.  At either level,
\begin{equation}
u_i=\tanh(Wh_i+b),\quad
\alpha_i=\frac{\exp(u_i^\top u_c)}{\sum_j\exp(u_j^\top u_c)},\quad
v=\sum_i\alpha_i h_i.
\end{equation}
The context vector $u_c$ learns which units matter for classification.  Hierarchy reduces a
variable document to semantically meaningful intermediate representations.

\begin{lstlisting}[caption={Reusable hierarchical attention pool}]
class AttnPool(nn.Module):
    def __init__(self, d):
        super().__init__(); self.proj = nn.Linear(d, d); self.context = nn.Parameter(torch.randn(d))
    def forward(self, h, mask=None):
        score = torch.tanh(self.proj(h)) @ self.context
        if mask is not None: score = score.masked_fill(~mask, -torch.inf)
        return torch.einsum('bt,btd->bd', score.softmax(-1), h)
\end{lstlisting}
\methodflow{words in sentences}{word encoder + attention}{sentence vectors}{sentence encoder + attention}{document class}{structurepurple}

\textbf{Evidence and limitations.} The original paper evaluates document classification on
Yelp, IMDB, Yahoo Answers, and Amazon review datasets against bag-of-words, convolutional,
and recurrent baselines using classification accuracy.  Word- and sentence-level attention
visualizations are useful diagnostics, not causal explanations.  The hierarchy is fixed at
two levels, pooling is lossy, bidirectional encoding is not a streaming solution, and the
model is designed for one document-level prediction rather than open-ended retrieval.

\textbf{PRA relation and experiment.} HAN supplies the principle that intermediate document
units deserve representations.  PRA generalizes those units to named sections, tables,
records, and code symbols and permits later expansion.  Compare flat chunks, HAN-style
sentence summaries, and URI hierarchy on classification plus evidence-required QA.  Remove
or substitute high-attention sentences and selected PRA nodes.  A strong hybrid uses HAN
pooling to construct cheap summaries while preserving children as resolvable objects.

\subsection{RAPTOR}

\textbf{Motivation and mechanism.} RAPTOR recursively embeds and clusters leaf chunks, asks
a language model to summarize each cluster, and repeats until a tree of increasingly abstract
nodes is formed \citep{sarthi2024raptor}.  With nodes $V_0$ at the leaves,
\begin{equation}
V_{\ell+1}=\{\operatorname{summarize}(C):
C\in\operatorname{cluster}(\{E(v):v\in V_\ell\})\}.
\end{equation}
Retrieval can search a collapsed set of all levels or traverse the tree, letting a query
choose between detailed evidence and broad synthesis.

\begin{lstlisting}[caption={Bottom-up RAPTOR index construction}]
level = chunk(document)
tree = [level]
while len(level) > 1:
    groups = cluster(embed(level))
    level = [summarizer(group) for group in groups]
    tree.append(level)
return tree
\end{lstlisting}
\methodflow{leaf text chunks}{embed + cluster}{recursive LLM summaries}{multi-level retrieval}{reader answer}{structurepurple}

\textbf{Evidence, paper trail, and failures.} RAPTOR reports QA results on QuALITY,
NarrativeQA, and QASPER, comparing conventional retrieval and long-context/retrieval
baselines with accuracy or task metrics.  The paper highlights a 20-point absolute QuALITY
gain for its GPT-4 configuration; that number is conditional on its model, tree, prompts,
and evaluation.  Descendants study dynamic maintenance and cheaper tree construction.
Indexing is expensive, abstractive summaries can omit or invent facts, clustering is unstable,
and an error at a high node can misroute an entire subtree.

\textbf{PRA relation and experiment.} RAPTOR chooses representations during indexing;
PRA chooses when and at which layer a representation becomes resident.  Assign immutable
URIs to every tree node and let the router escalate from summary to children.  Compare flat
RAG, collapsed-tree retrieval, explicit tree traversal, and PRA transport using the same
tree.  Measure answer quality, retrieved source tokens, LLM index cost, expansions, cache
reuse, summary corruption, and whether leaf evidence supports generated claims.

\subsection{MemWalker}

\textbf{Motivation and mechanism.} MemWalker recursively summarizes a long document into a
memory tree and has an LLM navigate it by choosing children until sufficient evidence is
found \citep{chen2023memwalker}.  At node $v_t$, a policy reads the query and child summaries:
\begin{equation}
a_t\sim\pi_\theta(a\mid q,s(v_t),\{s(c):c\in\mathrm{child}(v_t)\}),
\qquad v_{t+1}=\mathrm{child}(v_t,a_t).
\end{equation}
The traversal trace externalizes intermediate access decisions instead of retrieving all
nodes in one ranking call.

\begin{lstlisting}[caption={Bounded tree walk}]
node = root
for depth in range(max_depth):
    action = navigator.choose(question, node.summary,
                              [c.summary for c in node.children])
    if action == "answer": break
    node = node.children[action]
return reader(question, node.source_text)
\end{lstlisting}
\methodflow{question + root}{inspect child summaries}{LLM chooses branch}{repeat tree walk}{read leaf + answer}{structurepurple}

\textbf{Evidence and failure modes.} MemWalker evaluates long-context reasoning/QA against
retrieval and long-context baselines, measuring answer quality and traversal behavior.  It
is closer to an agent than static RAG, but the action space is confined to a prebuilt tree.
Serial LLM calls add latency, greedy early choices can make relevant siblings unreachable,
summary errors compound, and navigation traces show decisions without proving causal use
of the selected leaf.

\textbf{PRA relation and experiment.} Both implement progressive disclosure.  MemWalker
serializes branch decisions as language-model actions; PRA proposes a latent router inside
selected layers.  Use the same tree and compare textual navigation, latent top-$k$ child
selection, oracle paths, and a hybrid that escalates ambiguous nodes to the agent.  Match
node-expansion budgets and report latency, depth, backtracking, path recall, valid/shuffled
leaf effects, and calibration at the escalation boundary.

\subsection{PageIndex}

\textbf{Motivation and mechanism.} PageIndex builds a table-of-contents-like tree from
natural document sections and uses an LLM to reason over titles, summaries, and page ranges
rather than relying on vector similarity \citep{pageindex2026}.  A typical node is
$(\mathrm{id},\mathrm{title},[p_s,p_e],\mathrm{summary},\mathrm{children})$; retrieval is a
reasoned tree search followed by page/section extraction.

\begin{lstlisting}[caption={Reasoning-first section traversal}]
frontier = [root]
while budget.remaining() and not evidence_sufficient(frontier, query):
    decision = llm_rank(query, describe(frontier))
    frontier = expand_selected(frontier, decision.node_ids)
pages = merge_ranges([n.page_range for n in frontier])
return answer(query, pdf.extract(pages))
\end{lstlisting}
\methodflow{PDF/Markdown structure}{LLM-built section tree}{reason over node descriptions}{drill into page ranges}{answer with trace}{structurepurple}

\textbf{Evidence status and reproducibility.} PageIndex is an evolving open-source/product
system, not a settled peer-reviewed architecture result.  This survey pins its inspection
target to repository commit
\href{https://github.com/VectifyAI/PageIndex/commit/7592163e2a376b3917181fff9ac1858dc5daa2c6}{\texttt{7592163e2a37}}
(mirror snapshot synchronized 15 June 2026) and records the repository's then-documented
default \texttt{gpt-4o-2024-11-20}; access date is 3 August 2026.  The repository advertises
98.7\% FinanceBench accuracy, but this remains a vendor-reported compound-system result.
Reproduction must freeze PDFs/questions, parser/OCR, tree-building and answering prompts,
model versions, judge, exact-match policy, retries, and dollar/token cost.  Natural structure
can itself be missing or malformed, and every LLM tree call can introduce nondeterminism.

\textbf{PRA relation and experiment.} PageIndex supplies a strong logical namespace for
PRA: node IDs/page ranges become versioned URIs, and selected sections can materialize as
text or per-layer KV.  Compare vector RAG, PageIndex traversal, PRA latent traversal, and
agent escalation on the same tree.  Stratify exact lookup, cross-reference, multi-section
synthesis, and cross-document reasoning.  Report source-page recall, answer/citation quality,
tree calls, latency, cost, and invalidation after a PDF revision.

\subsection{GraphRAG}

\textbf{Motivation and mechanism.} GraphRAG extracts entities, relations, and claims from a
corpus, clusters the resulting graph into communities, and creates community reports
\citep{edge2024graphrag}.  Local search expands around query entities; global search maps a
question over community reports and reduces partial answers:
\begin{equation}
\hat y=\operatorname{reduce}\{\operatorname{answer}(q,R_c):
c\in\operatorname{select}(q,\mathcal C)\}.
\end{equation}
The graph exposes corpus-wide themes and multi-hop relationships that may not be similar to
the surface query.

\begin{lstlisting}[caption={GraphRAG global map-reduce sketch}]
reports = select_community_reports(query, graph.communities, token_budget)
partials = [llm.map_answer(query, report) for report in reports]
return llm.reduce_answer(query, partials)
\end{lstlisting}
\methodflow{source corpus}{extract entity/relation graph}{community summaries}{local or global search}{map-reduce answer}{structurepurple}

\textbf{Evidence and failure modes.} The GraphRAG preprint studies global sensemaking over
roughly million-token corpora and reports LLM-judged comprehensiveness and diversity relative
to conventional RAG.  These are inference-time judgments over an expensive LLM-derived
index, so model, prompts, judge, and cost are integral conditions.  Extraction errors,
entity resolution, hallucinated relations, index cost, stale derived summaries, and judge
bias are prominent risks.  Local factual QA and global thematic synthesis are distinct tasks.

\textbf{PRA relation and experiment.} A PRA URI can address a node, neighborhood, graph
query result, or community report.  The hard problem becomes dependency-aware invalidation:
changing one source may stale many derived KVs.  Compare flat RAG, GraphRAG text transport,
and graph-node PRA transport on local and global questions.  Match source and answer models;
report index cost separately from per-query cost, provenance to original spans, stale-update
tests, and counterfactual removal of decisive nodes.

\subsection{PaperTrail as provenance and HCI}

\textbf{Motivation and mechanism.} PaperTrail is not a context architecture.  It decomposes
an answer and its sources into claims and evidence, then visualizes supported claims,
unsupported assertions, and omitted source information \citep{martinboyle2026papertrail}.
If $C$ is answer claims and $E$ source evidence units, the interface exposes a relation
$R\subseteq C\times E$ plus unsupported $C\setminus\operatorname{dom}(R)$ and uncovered
source claims.

\begin{lstlisting}[caption={Claim--evidence audit pipeline}]
answer_claims = decompose_claims(answer)
source_claims = decompose_claims(sources)
links = align_and_verify(answer_claims, source_claims)
render(answer_claims, source_claims, links,
       unsupported=set(answer_claims)-links.claims())
\end{lstlisting}
\methodflow{answer + sources}{claim decomposition}{claim--evidence alignment}{support/omission labels}{verification interface}{structurepurple}

\textbf{Evidence and limitations.} The 2026 work evaluates PaperTrail in a within-subjects
study with 26 researchers performing scholarly editing tasks.  The interface lowered trust
relative to a baseline but did not produce corresponding behavioral change: participants
still relied on generated edits when verification was cognitively costly.  This is an HCI
result about provenance presentation, not evidence for longer context or retrieval quality.
Claim decomposition and alignment can themselves be wrong, and more visible evidence can
overload rather than assist a user.

\textbf{PRA relation and experiment.} PRA's URI/version and routing traces can populate a
PaperTrail-like view, but attended memory is not causal support.  Record selected objects,
source spans, cache versions, generated claims, and outputs under removal/substitution.
Compare citation-only, route-trace, and counterfactual claim-evidence interfaces.  Measure
verification accuracy, time, calibrated trust, unsupported-claim detection, omissions, and
cognitive load.  This positions PaperTrail correctly as provenance/HCI rather than a
competitor in the context architecture table.

\section{Category VII: Agents with Search and Tool Use}

Agents move relevance decisions outside a single forward pass.  They can discover unknown
sources, reformulate a question after observing evidence, act on the world, and backtrack.
Their generality costs serial latency, tokens, tool failures, and a larger audit surface.
PRA should be compared as a restricted fast path for repeated read-only access, not as a
replacement for open-ended action.

\subsection{ReAct}

\textbf{Motivation and mechanism.} ReAct interleaves natural-language reasoning with
actions and observations, allowing evidence acquisition to change the next reasoning step
\citep{yao2023react}.  A trajectory is
\begin{equation}
\tau=(q,r_1,a_1,o_1,r_2,a_2,o_2,\ldots,y),\qquad
(r_t,a_t)\sim\pi_\theta(\cdot\mid q,\tau_{<t}).
\end{equation}
Reasoning maintains a task state; tools ground or alter that state; the policy decides when
to stop.  Search is only one environment---the paper also studies interactive tasks.

\begin{lstlisting}[caption={ReAct execution loop with an explicit budget}]
trace = []
for step in range(max_steps):
    thought, action = model.next(question, trace)
    if action.name == "finish": return action.answer, trace
    observation = tools[action.name](**action.arguments)
    trace.append((thought, action, observation))
raise BudgetExceeded(trace)
\end{lstlisting}
\methodflow{question + trace}{reason}{tool action}{observation}{repeat or finish}{agentred}

\textbf{Evidence, descendants, and failures.} ReAct evaluates HotpotQA and FEVER plus
ALFWorld and WebShop, comparing standard prompting, chain-of-thought, act-only, and ReAct
using exact match, factual accuracy, or environment success.  Prompt examples, base model,
tool interface, and action budget materially condition results.  Descendants add planning,
reflection, memory, and richer tools.  Common failures are reasoning loops, malformed calls,
irrelevant searches, premature stopping, prompt injection in observations, and high latency.

\textbf{PRA relation and experiment.} Both progressively acquire context.  ReAct remains
necessary when sources are unknown, actions have side effects, or explicit planning and
backtracking are essential.  Compare ReAct search, PRA over a known reference graph, and a
tiered policy that tries local context, cached URI memory, structural navigation, then search.
Match evidence/action budgets and report answer quality, calls, tokens, latency, provenance,
cache reuse, and escalation calibration.  Cache successful read-only discoveries as versioned
references, never silently replay side-effecting actions.

\subsection{Self-Ask with Search}

\textbf{Motivation and mechanism.} Self-Ask addresses compositionality by explicitly
generating follow-up questions.  An external search component answers each subquestion, and
the model combines intermediate answers into the final response
\citep{press2022selfask}:
\begin{equation}
q_t=\operatorname{decompose}(q,a_{<t}),\quad
e_t=\operatorname{search}(q_t),\quad
y=\operatorname{compose}(q,\{q_t,e_t\}).
\end{equation}
The textual protocol makes the retrieval chain easy to inspect and separates decomposition
errors from search errors.

\begin{lstlisting}[caption={Self-Ask protocol}]
facts = []
while model.needs_followup(question, facts):
    subquestion = model.followup(question, facts)
    evidence = search(subquestion)
    facts.append((subquestion, evidence))
return model.final_answer(question, facts)
\end{lstlisting}
\methodflow{compositional question}{generate follow-up}{search intermediate answer}{accumulate facts}{compose final answer}{agentred}

\textbf{Evidence and limitations.} The work studies multi-hop/compositional question
answering and the ``compositionality gap'', comparing direct answer, chain-of-thought, and
search-assisted decomposition with exact match.  Decomposition improves observability, but
one wrong subquestion can derail all later hops; intermediate search snippets may be noisy;
serial calls cost latency; and a plausible intermediate answer can conceal missing evidence.

\textbf{PRA relation and experiment.} PRA may represent a follow-up target as a generated
reference only when a resolver can bind it deterministically; free-form web search is not a
URI lookup.  Compare text subquestions, latent reference queries, and a hybrid that uses
PRA for known namespaces and search otherwise.  Use the same search backend and hop budget;
measure decomposition accuracy, evidence recall, calls, latency, valid/shuffled reference
effects, and failures where the correct source is absent from the advertised graph.

\subsection{IRCoT}

\textbf{Motivation and mechanism.} IRCoT interleaves retrieval with chain-of-thought so
partial reasoning becomes a better query for the next evidence hop
\citep{trivedi2023ircot}.  Rather than retrieve once from the original question,
\begin{equation}
z_t=R(q,r_{<t},e_{<t}),\qquad
r_t=G(q,r_{<t},e_{<t},z_t),
\end{equation}
and the cycle continues until an answer or budget limit.  Retrieval and reasoning therefore
co-evolve.

\begin{lstlisting}[caption={IRCoT-style retrieval/reasoning alternation}]
state, evidence = question, []
for hop in range(max_hops):
    passages = retriever.search(state, k)
    evidence.extend(passages)
    state, answer = reasoner.step(question, state, passages)
    if answer is not None: return answer, evidence
\end{lstlisting}
\methodflow{question + partial rationale}{retrieve passages}{extend evidence}{reason next hop}{answer or retrieve again}{agentred}

\textbf{Evidence, paper trail, and failures.} IRCoT evaluates multi-step QA including
HotpotQA, 2WikiMultiHopQA, MuSiQue, and related datasets, comparing retrieve-then-read,
iterative retrieval, and chain-of-thought configurations with exact match/F1 and retrieval
metrics.  It motivates later iterative and agentic RAG.  Reasoning text can introduce false
entities that poison later queries, cost grows by hop, stopping is difficult, and answer
gains can reflect a stronger prompt rather than a better retriever.

\textbf{PRA relation and experiment.} PRA layers or decoding states could emit successive
latent queries, approximating IRCoT without serializing each query.  Compare textual query
rewriting and latent queries against the same frozen retriever, plus oracle evidence order.
Match retrieval calls/candidates and report answer quality, hop recall, latency, and
counterfactual evidence effects.  The hybrid should expose an auditable object path even if
queries are latent; otherwise reduced tokens come at the cost of provenance.

\subsection{Toolformer}

\textbf{Motivation and mechanism.} Toolformer teaches a language model when and how to call
APIs by sampling candidate calls and retaining those that reduce language-model loss
\citep{schick2023toolformer}.  For call $c$ inserted at position $i$, a simplified filter is
\begin{equation}
\Delta(c)=L(x_{i:}\mid x_{<i})-L(x_{i:}\mid x_{<i},c,\operatorname{result}(c));
\quad c\text{ kept if }\Delta(c)>\tau.
\end{equation}
Fine-tuning on retained demonstrations yields discrete call triggering and result integration
without dense human call labels.

\begin{lstlisting}[caption={Loss-improvement filtering for pseudo-labeled calls}]
for position in candidate_positions(text):
    for call in sample_calls(model, text, position):
        result = execute_sandboxed(call)
        gain = loss_without(call) - loss_with(call, result)
        if gain > threshold: training_examples.append(insert(text, call, result))
\end{lstlisting}
\methodflow{raw text}{sample API calls}{execute + score loss gain}{self-supervised call data}{tool-aware LM}{agentred}

\textbf{Evidence and limitations.} Toolformer evaluates a 6.7B GPT-J model with calculator,
QA, search, translation, and calendar tools on zero/few-shot downstream tasks and language
modeling, comparing base and tool-aware models.  Tool choice and API formatting ablations
matter.  Pseudo-labels inherit the base model's blind spots, unsafe or costly calls require
sandboxing, tool errors enter the context, and lower token loss does not guarantee factual
or task utility.

\textbf{PRA relation and experiment.} Toolformer's call-retention principle is directly
useful for PRA router training: sample candidate object accesses, keep those whose controlled
insertion lowers loss, then train a gate with corrupted negatives.  Compare heuristic cosine,
supervised relevance, and loss-filtered routers under equal trainable parameters.  Evaluate
selection recall, calibration, task loss, valid/shuffled/irrelevant effects, and access cost.
PRA can internalize deterministic reads; arbitrary or side-effecting APIs should remain
explicit tool calls.

\subsection{WebGPT and browser agents}

\textbf{Motivation and mechanism.} WebGPT trains a model to issue search queries, open
results, navigate pages, collect quotations, and answer with citations, using demonstrations,
behavior cloning, reward modeling, and human feedback \citep{nakano2021webgpt}.  The browser
state makes information acquisition a sequential decision process:
\begin{equation}
a_t\sim\pi_\theta(a\mid s_t),\quad s_{t+1}=\operatorname{browser}(s_t,a_t),
\quad y=G(q,\operatorname{quotes}(\tau)).
\end{equation}

\begin{lstlisting}[caption={Citation-preserving browser loop}]
state, evidence = browser.start(question), []
for _ in range(action_budget):
    action = policy(state, evidence)
    if action.kind == "quote": evidence.append(browser.quote(action.span))
    elif action.kind == "answer": return compose_with_citations(question, evidence)
    else: state = browser.execute(action)
\end{lstlisting}
\methodflow{question}{search/open/navigate}{quote source spans}{human-feedback policy}{cited answer}{agentred}

\textbf{Evidence, descendants, and failures.} WebGPT evaluates long-form question answering
using human preference, factuality, and citation-related judgments against retrieval and
human-written baselines.  It established a lineage toward browser-assisted research agents.
Results are bound to the web snapshot, search engine, model, demonstrations, reward model,
and evaluator.  Failures include source-quality errors, selective quotation, citation drift,
prompt injection, browsing loops, changing pages, and high serial latency.

\textbf{PRA relation and experiment.} PRA cannot discover an unknown open-web source.
It can cache already discovered pages, versions, sections, and evidence for later questions.
Compare fresh browser search, cached text RAG, cached PRA KVs, and a policy that searches on
cache miss or staleness.  Measure answer/citation correctness, source diversity, actions,
latency, cache hit and avoided encoding, update detection, and security under malicious page
instructions.  Provenance must retain the original page/version even when KVs are reused.

\subsection{Recursive Language Models and context as environment}

\textbf{Motivation and mechanism.} Recursive language-model approaches treat an oversized
context as an external environment.  A model writes code or issues subcalls that inspect,
partition, summarize, and recursively solve pieces instead of consuming the entire source
in one forward pass.  For context $D$ and budget $B$,
\begin{equation}
y=F(q,D,B)=\operatorname{aggregate}\{F(q_j,D_j,B_j)\}_{j=1}^m,
\qquad \sum_j B_j\leq B.
\end{equation}
The model controls the decomposition and can adapt it to each question.

\begin{lstlisting}[caption={Budgeted recursive context program}]
def solve(query, view, budget):
    if view.fits() or budget.depth == 0:
        return model(query, view.read(budget.tokens))
    plan = model.plan_partitions(query, view.metadata(), budget)
    partials = [solve(p.subquery, view.slice(p.range), budget.child(p)) for p in plan]
    return model.aggregate(query, partials)
\end{lstlisting}
\methodflow{oversized context environment}{plan partitions/search}{recursive model subcalls}{aggregate partial results}{final answer}{agentred}

\textbf{Evidence and failure modes.} The family is evaluated on inputs beyond a model's
window, code/document analysis, aggregation, and long-context reasoning, with answer quality,
subcall count, tokens, and cost.  Baselines should include truncation, long context, RAG, and
non-recursive map-reduce.  Generality is high, but recursive calls are expensive, partition
boundaries can destroy dependencies, summaries compound errors, arbitrary generated code
needs sandboxing, and budgets/stopping rules can dominate outcomes.

\textbf{PRA relation and experiment.} PRA can be viewed as a differentiable cached fast path
for recurring read operations discovered by an agent or recursive program.  Use identical
source partitions and compare RLM subcalls, tree navigation, PRA references, and a tiered
system that escalates on uncertainty.  Match total model FLOPs/tokens and wall-clock budget;
report depth, calls, source coverage, provenance, cache reuse, and tasks whose relevant
partition is unknown initially.  The agent should win on open-ended discovery; PRA must
justify itself through repeated low-latency reuse and causal content selection.

\Needspace{12\baselineskip}
\subsection{A tiered agent--PRA design}

The categories are best composed as a monotone escalation policy.  Local attention handles
dense immediate context; a warm URI cache handles repeated named reads; a tree/graph index
handles bounded structural search; and a general agent handles discovery or action.  Each
tier returns both evidence and calibrated confidence.  Expensive discoveries may be cached
only if the source is immutable or versioned, authorization permits reuse, and the action is
read-only.  The evaluation must include wrong-confidence costs: unnecessary escalation,
silent cache hits on stale data, and premature stopping before an agent search.

\methodflow{local context}{PRA URI cache}{tree/graph navigation}{external agent/tools}{answer + causal provenance}{praorange}

\iffalse
\subsection{Motivation and central insight}
PRA begins from the observation that practical language is highly referential. Documents refer to sections, tables, figures, records, code symbols, prior messages, and external resources. Flattening all referenced material into the prompt discards this structure and pays for all tokens before the model knows which ones matter. One-shot RAG retains an external store, but usually commits to a fixed set of retrieved passages before generation. Agentic search permits progressive disclosure, but at orchestration latency and with a coarse text/tool boundary.

PRA separates \emph{naming} from \emph{materialization}. A local sequence can contain a compact handle such as \texttt{doc://report/section/3}. The handle is a stable identity and provenance object. Its content may be resolved, encoded, cached, and exposed to selected layers only when a learned or runtime policy predicts utility. This turns retrieval into a context-access primitive near attention rather than only a preprocessing stage.

\subsection{Standalone prototype architecture}
For layer $\ell$, ordinary causal attention is
\begin{equation}
A^{(\ell)}_{\mathrm{local}}=\operatorname{softmax}\left(\frac{Q^{(\ell)}K^{(\ell)\top}}{\sqrt{d_h}}+\mathcal{M}_{\mathrm{causal}}\right)V^{(\ell)}.
\end{equation}
Each candidate reference $r$ is resolved into content $z_r$ and encoded into layer-specific memory $(K_r^{(\ell)},V_r^{(\ell)})$. A router scores summaries $s_r^{(\ell)}$:
\begin{equation}
\alpha_r^{(\ell)} = \frac{\langle \bar q^{(\ell)},s_r^{(\ell)}\rangle}{\sqrt{d_h}}, \qquad \mathcal{C}^{(\ell)}=\operatorname{TopK}_r \alpha_r^{(\ell)}.
\end{equation}
The memory branch is
\begin{equation}
A^{(\ell)}_{\mathrm{mem}}=\operatorname{softmax}\left(\frac{Q_m^{(\ell)}K_{\mathcal{C}}^{(\ell)\top}}{\sqrt{d_h}}+\mathcal{M}_{\mathrm{ref}}\right)V_{\mathcal{C}}^{(\ell)},
\end{equation}
and the residual update is
\begin{equation}
H^{(\ell+1)}=H^{(\ell)}+A^{(\ell)}_{\mathrm{local}}+g^{(\ell)}A^{(\ell)}_{\mathrm{mem}}+\operatorname{MLP}^{(\ell)}(\cdot),
\end{equation}
with a gate initialized near zero to preserve pretrained behavior.

\begin{figure}[h]
\centering
\begin{tikzpicture}[node distance=8mm and 10mm, every node/.style={font=\small}, box/.style={draw,rounded corners,align=center,minimum height=8mm,fill=blue!4}, arr/.style={-{Latex[length=2mm]}}]
\node[box] (tok) {local tokens +\\reference handles};
\node[box,right=of tok] (sa) {causal\\self-attention};
\node[box,below=of tok] (res) {URI resolver\\and object store};
\node[box,right=of res] (enc) {per-layer encoder\\and KV cache};
\node[box,right=of enc] (route) {query-dependent\\router / top-$k$};
\node[box,right=of sa] (fuse) {gated fusion};
\node[box,right=of fuse] (out) {next hidden\\state};
\draw[arr] (tok)--(sa); \draw[arr] (tok)--(res); \draw[arr] (res)--(enc); \draw[arr] (enc)--(route); \draw[arr] (sa)--(fuse); \draw[arr] (route)--(fuse); \draw[arr] (fuse)--(out);
\end{tikzpicture}
\caption{PRA separates lightweight reference naming from latent memory materialization.}
\end{figure}

\subsection{Pseudocode and PyTorch sketch}
\begin{lstlisting}[caption={PRA layer pseudocode}]
for layer_id, layer in enumerate(layers):
    h_local = layer.self_attention(h, causal_mask)
    refs = parse_reference_handles(input_ids, metadata)
    candidates = cache.lookup_or_encode(refs, layer_id, resolver)
    selected = router.topk(layer.query_summary(h), candidates)
    h_mem = layer.cross_attention(h, selected.keys, selected.values)
    h = h + h_local + sigmoid(layer.gate_logit) * h_mem
    h = h + layer.mlp(layer.norm(h))
\end{lstlisting}
\begin{lstlisting}[caption={Minimal educational PRA module}]
class PRABlock(nn.Module):
    def __init__(self, d_model, n_heads):
        super().__init__()
        self.local = nn.MultiheadAttention(d_model, n_heads, batch_first=True)
        self.memory = nn.MultiheadAttention(d_model, n_heads, batch_first=True)
        self.gate_logit = nn.Parameter(torch.tensor(-4.0))
        self.norm = nn.LayerNorm(d_model)
        self.ff = nn.Sequential(nn.Linear(d_model, 4*d_model), nn.GELU(),
                                nn.Linear(4*d_model, d_model))

    def forward(self, x, mem, causal_mask=None, mem_pad_mask=None):
        y, _ = self.local(x, x, x, attn_mask=causal_mask, need_weights=False)
        if mem is not None and mem.size(1):
            m, _ = self.memory(x, mem, mem,
                               key_padding_mask=mem_pad_mask,
                               need_weights=False)
            y = y + torch.sigmoid(self.gate_logit) * m
        x = x + y
        return x + self.ff(self.norm(x))
\end{lstlisting}

\subsection{Preliminary experiments in Paper 1}
The prototype compares four-layer width-256 SelfAttention, four-layer PRA, and a two-SelfAttention/two-PRA hybrid on WikiText-2, with a 2,000-token BPE vocabulary, context length 128, seeds 1 and 7, and 524,288 processed tokens per run. Mean plain-text test losses were 4.7691, 4.7457, and 4.7574. The defensible conclusion is functional equivalence at this pilot scale, not superiority, because same-width models were not parameter matched.

In a second phase, entries were split into one to five URI-addressed earlier parts and a held-out natural-language tail. The PRA base was frozen and zero-initialized memory output projections were trained. Valid references reduced answer-token loss relative to a disabled path by 0.1165 (full PRA) and 0.0841 (hybrid). Cross-example shuffled references increased loss by only 0.0145 and 0.0122. Thus the branch contributes, and there is preliminary content sensitivity, but the smaller shuffled delta warns that much of the benefit could be generic adaptation or weakly discriminative memory use.

\subsection{Evaluation doctrine}
PRA's strongest methodological contribution may be its insistence on counterfactual reference evaluation. A task score alone does not establish retrieval use. Required conditions include valid, disabled, shuffled, irrelevant, empty, random-ID, content-only, position-only, oracle, and all-content-local controls. The principal causal quantities are
\begin{equation}
\Delta_{\mathrm{use}}=L_{\mathrm{disabled}}-L_{\mathrm{valid}},\qquad
\Delta_{\mathrm{content}}=L_{\mathrm{shuffled}}-L_{\mathrm{valid}}.
\end{equation}
A credible result requires $\Delta_{\mathrm{content}}>0$, increasing advantage on locally insufficient examples, and routing that ranks known relevant parts above distractors.

\subsection{Paper trail and open claims}
Paper 0 develops the architecture and systems vision; Paper 1 turns it into an inspectable PyTorch instrument and narrows the empirical claims. The next lineage should include parameter-matched controls, natural QA and long-context benchmarks, cache-reuse measurements, selective triggering, pretrained-model adaptation, and direct comparisons against long context, RAG, latent retrieval, structural indexes, and agents. The central unproven thesis is efficiency through reusable, demand-driven latent expansion.

\section{Category I: Dense, Sparse, and Approximate Long-Context Attention}
\subsection{Full attention and positional extension}
Full self-attention makes every token visible to every compatible token. RoPE scaling, interpolation, and long-context continued pretraining can expand the usable positional range without changing the dense $O(T^2)$ graph. The approach is maximally general and hardware-friendly at moderate lengths, but pays for irrelevant tokens and offers no persistent object identity.

\textbf{Experiment pattern.} Needle-in-a-haystack, passkey retrieval, long-document QA, perplexity by position, and throughput/memory versus context length. For PRA, these become the ``all evidence locally visible'' upper control and an equal-memory/equal-compute baseline.

\subsection{Longformer}
Longformer combines a sliding local window with selected global tokens, reducing attention complexity from quadratic to approximately $O(Tw+Tg)$ for window $w$ and $g$ global positions \citep{beltagy2020longformer}. Its insight is that most dependencies are local while a small number of task tokens require global connectivity.
\begin{equation}
\mathcal{N}(i)=\{j:|i-j|\leq w/2\}\cup \mathcal{G}.
\end{equation}
\begin{lstlisting}
# educational local+global mask
idx = torch.arange(T)
local = (idx[:, None] - idx[None, :]).abs() <= window // 2
global_mask = global_tokens[:, None] | global_tokens[None, :]
allowed = local | global_mask
scores = scores.masked_fill(~allowed, float('-inf'))
\end{lstlisting}
Original experiments included character-level language modeling and long-document classification/QA, followed by LED for long-input generation. Longformer influenced domain models such as Clinical-Longformer. Relative to PRA, it makes relevance a static token-graph decision rather than named object selection. A useful comparison fixes total source tokens and varies sparsity of relevant evidence: Longformer should excel when global positions are known; PRA should benefit when a compact handle can defer a large object.

\subsection{BigBird}
BigBird mixes local, random, and global edges, achieving linear sparse attention while retaining theoretical universal approximation and Turing-completeness properties under its graph assumptions \citep{zaheer2020bigbird}. Its original evaluations covered QA, summarization, and genomics, including sequences substantially longer than standard BERT.
\begin{equation}
\mathcal{N}(i)=\mathcal{N}_{\mathrm{local}}(i)\cup\mathcal{N}_{\mathrm{random}}(i)\cup\mathcal{G}.
\end{equation}
BigBird's random edges provide rapid graph mixing, while PRA's router provides content-dependent edges to named external memories. The clean experiment is a matched sparse-edge budget: random/global token edges versus top-$k$ reference memories, evaluated under relevant-reference corruption.

\subsection{Performer and kernelized attention}
Performer approximates softmax attention using positive random features, rewriting attention as associative products and reducing sequence complexity \citep{choromanski2021performer}. If $\exp(q^\top k)\approx \phi(q)^\top\phi(k)$,
\begin{equation}
\operatorname{Attn}(Q,K,V)\approx D^{-1}\phi(Q)\left(\phi(K)^\top V\right).
\end{equation}
This is a computation approximation, not a relevance architecture: all tokens still contribute through a compressed statistic. PRA can use linear attention locally or inside memory encoding, but preserves discrete identity and counterfactual replacement. Compare approximation error, exact retrieval, and memory reuse separately.

\section{Category II: Recurrence and Compression}
\subsection{Transformer-XL}
Transformer-XL introduces segment-level recurrence by reusing previous hidden states as stop-gradient memory and uses relative positions to avoid temporal confusion \citep{dai2019transformerxl}. For segment $s$,
\begin{equation}
\tilde H_{s-1}^{(\ell)}=\operatorname{SG}(H_{s-1}^{(\ell)}),\quad K,V=[\tilde H_{s-1}^{(\ell)};H_s^{(\ell)}].
\end{equation}
Its experiments showed improved language modeling and longer effective dependencies. The paper trail includes XLNet and many recurrent-memory variants. Transformer-XL memory is chronological and implicit; PRA memory is addressable and can be non-contiguous. A hybrid could use recurrent local memory for recent discourse and PRA for remote named objects.

\subsection{Compressive Transformer}
Compressive Transformer retains a short-term memory and compresses evicted activations into a lower-resolution long-term memory \citep{rae2020compressive}. If $C$ is a compression function,
\begin{equation}
M^{c}_{s}= [M^{c}_{s-1}; C(M^{m}_{s-1,\mathrm{evicted}})].
\end{equation}
The key trade-off is rate versus retained predictive information. Experiments examined language modeling and which compression objectives preserve long-range information. PRA differs by retrieving original or cached object representations rather than obligatorily compressing chronological history. PRA could attach multiple materializations to one URI: summary, compressed latent, and full KV, escalating by uncertainty.

\subsection{Recurrent Memory Transformer (RMT)}
RMT places learned memory tokens at segment boundaries. Their output representations are carried into the next segment, allowing end-to-end learned compression \citep{bulatov2022rmt}. It was evaluated on language modeling, algorithmic tasks, and later million-token synthetic tasks. RMT asks the model to decide what to preserve during reading; PRA asks what to fetch during use. Combining them gives a write-time summary plus read-time dereferencing.

\subsection{Infini-attention}
Infini-attention combines local causal attention with a compressive associative memory updated across segments \citep{munkhdalai2024infini}. A linear memory can be updated as
\begin{equation}
M_s \leftarrow M_{s-1}+\phi(K_s)^\top V_s,\qquad z_s\leftarrow z_{s-1}+\sum_i\phi(k_i),
\end{equation}
then queried using $\phi(Q_s)M_s/(\phi(Q_s)z_s)$. Experiments targeted long-context language modeling, passkey retrieval, and million-token summarization. Infini-attention offers bounded memory but lossy superposition; PRA offers selective object identity but potentially unbounded external storage. Evaluate exact-value recall under collisions and semantic synthesis under compression.

\section{Category III: Differentiable External Memory}
\subsection{Memory Networks and Differentiable Neural Computers}
Memory Networks separate input encoding, memory storage, multi-hop addressing, and response generation \citep{weston2015memory}; DNC adds differentiable read/write heads, content lookup, temporal links, and allocation \citep{graves2016dnc}. A standard content read is
\begin{equation}
w_i=\operatorname{softmax}(\beta\,\operatorname{sim}(q,m_i)),\qquad r=\sum_i w_i m_i.
\end{equation}
These systems established learned multi-hop memory access on synthetic QA and algorithmic tasks. Their weakness at modern scale is complex write dynamics and optimization. PRA narrows the problem: object writes can be handled by a resolver/cache, while learned attention focuses on read selection.

\subsection{Memformer}
Memformer uses a fixed number of dynamic memory slots and memory replay backpropagation to process long sequences with constant memory \citep{wu2020memformer}. Experiments reported competitive performance with lower memory and faster inference. Its slots are learned summaries without stable semantic addresses. PRA could use Memformer-like slots as an object-level cache tier while preserving URI provenance in metadata.

\subsection{Memorizing Transformer}
Memorizing Transformer augments a Transformer with a non-differentiable kNN memory of past key/value pairs, querying a large episodic store during language modeling \citep{wu2022memorizing}. For query $q$:
\begin{equation}
\mathcal{I}=\operatorname{kNN}(q,\{k_i\}),\quad a=\operatorname{softmax}(qK_{\mathcal I}^\top)V_{\mathcal I}.
\end{equation}
The paper reported perplexity gains across datasets and scales up to billions of parameters, with benefits increasing with memory size before diminishing returns. It strongly anticipates PRA's latent memory branch. The distinction is explicit reference identity and resolvable object boundaries versus similarity over an episodic stream. Compare open kNN over all prior states against two-stage URI candidate restriction plus within-object attention.

\section{Category IV: Retrieval-Augmented Pretraining and Generation}
\subsection{kNN-LM}
kNN-LM stores training-context representations and next-token labels. At inference, a nearest-neighbor distribution is interpolated with the LM distribution \citep{khandelwal2020knnlm}:
\begin{equation}
p(y|q)=\lambda p_{\mathrm{kNN}}(y|q)+(1-\lambda)p_{\mathrm{LM}}(y|q).
\end{equation}
Its experiments showed perplexity improvements and domain adaptation without model retraining. PRA instead retrieves rich memories that influence hidden computation, not only the final token distribution. A direct comparison should measure calibration and compositional use when retrieved evidence is not a near-duplicate continuation.

\subsection{REALM}
REALM jointly pretrains a retriever and masked language model over a large document corpus, marginalizing retrieved documents approximately \citep{guu2020realm}. It demonstrated open-domain QA gains and retrieval of world knowledge. The retriever is trained from the language objective, a principle PRA should adopt through auxiliary reference supervision or end-to-end marginalization. REALM materializes retrieved text in an encoder; PRA can cache latent per-layer representations and resolve explicit references.

\subsection{RAG}
RAG combines a dense retriever with a seq2seq generator and marginalizes over documents either per sequence or per token \citep{lewis2020rag}:
\begin{equation}
p(y|x)=\sum_{z\in\operatorname{topk}(x)}p_\eta(z|x)\prod_t p_\theta(y_t|x,z,y_{<t}).
\end{equation}
Original experiments covered open-domain QA and knowledge-intensive generation. RAG became a dominant systems pattern and inspired reranking, query rewriting, corrective retrieval, and agentic RAG. PRA's strongest contrast is when and how retrieval enters: text passages selected before generation versus addressable latent memories available by layer and step. Fair comparisons must include equivalent source corpora, retrieval budgets, and citations/provenance.

\subsection{RETRO}
RETRO retrieves nearest-neighbor chunks from a massive token database and interleaves chunked cross-attention in a decoder \citep{borgeaud2022retro}. It showed that retrieval from trillions of tokens can match substantially larger parametric models. RETRO is one of PRA's closest ancestors: both use a dedicated cross-attention path to external context. PRA adds explicit handles, potentially reusable per-object KV caches, progressive triggering, and a stronger naming/materialization distinction. A decisive comparison tests no explicit handle, noisy handle, exact handle, and learned global retrieval under the same corpus.

\section{Category V: Attention-Level and KV Retrieval}
\subsection{Unlimiformer}
Unlimiformer wraps pretrained encoder-decoder models by indexing encoder hidden states and using kNN search to compute exact top-$k$ cross-attention candidates \citep{bertsch2023unlimiformer}. It was evaluated on long-document summarization and related seq2seq tasks. Its important insight is that maximum inner-product search can serve as sparse attention selection without retraining the whole model.

PRA and Unlimiformer converge at latent retrieval. Unlimiformer indexes token states from the current long input; PRA indexes named external objects and may persist their layer-specific KVs across requests. An adaptation experiment should compare (a) full cross-attention, (b) Unlimiformer token MIPS, (c) PRA object routing then token attention, and (d) hierarchical object-to-token retrieval.

\subsection{PagedAttention and prefix caching}
PagedAttention is a serving-memory architecture rather than a semantic retrieval model. vLLM partitions KV caches into blocks, reducing fragmentation and enabling sharing \citep{kwon2023vllm}. Prefix caching reuses identical prefix blocks across requests. This work changes the feasibility of PRA: URI-addressed objects can map naturally to immutable KV block sets with reference counts, eviction, isolation, and copy-on-write.

The comparison is therefore not competitive but compositional. PRA defines \emph{which external memories are semantically addressable}; a PagedAttention-like runtime defines \emph{how those tensors are allocated and shared}. Required systems experiments include cache hit rate, bytes reused, resolve latency, encode latency, synchronized end-to-end latency, fragmentation, and multi-tenant isolation.

\subsection{KV compression and token eviction}
StreamingLLM, H$_2$O, Scissorhands, SnapKV, and related systems retain a subset of past KV entries using attention sinks, accumulated importance, or query-aware selection. They optimize the active cache of a single stream rather than retrieving an external knowledge object. PRA could use these methods inside each object cache and expose multiple fidelity levels. The key evaluation difference is that eviction methods should be judged by retention of necessary past states, whereas PRA must additionally prove correct object selection and content causality.

\section{Category VI: Structural and Hierarchical Indexing}
\subsection{Hierarchical Attention Networks}
Hierarchical Attention Networks encode words into sentence vectors and sentences into document vectors with attention at both levels \citep{yang2016han}. Although predating modern RAG, HAN supplies a crucial principle: document structure creates meaningful intermediate units. PRA can generalize this from fixed word/sentence levels to URI-addressed sections, tables, paragraphs, records, and code symbols.

\subsection{RAPTOR}
RAPTOR recursively embeds, clusters, and summarizes chunks, building a tree with multiple abstraction levels \citep{sarthi2024raptor}. At query time, retrieval can select leaf text or higher-level summaries. Its experiments showed improvements on complex QA, including a large absolute gain on QuALITY when combined with GPT-4. The paper trail includes dynamic RAPTOR variants and cheaper tree-construction methods.

RAPTOR decides materializations during indexing; PRA decides whether and where to materialize during inference. A natural synthesis assigns URIs to every tree node. Upper PRA layers query abstract summaries; lower or later passes expand selected children. Experiments should measure answer quality, retrieved tokens, tree expansions, and robustness to summary errors.

\subsection{MemWalker}
MemWalker has an LLM walk through a hierarchy of recursively generated summaries for long-document reasoning \citep{chen2023memwalker}. It transforms retrieval into sequential navigation: inspect a node, choose a child, and continue. This is closer to an agent than static RAG, but its action space is bounded by a memory tree. PRA could internalize the child-selection step as a layer router while retaining explicit traversal traces.

\subsection{PageIndex}
PageIndex is a vectorless, reasoning-first document retrieval system that builds a table-of-contents-like tree and lets an LLM reason over section descriptions and drill down. Its public implementation emphasizes traceability and natural document structure rather than fixed chunks or vector similarity \citep{pageindex2026}. Because its claims are presently more system- and benchmark-driven than established peer-reviewed architecture science, comparisons should reproduce datasets, judges, prompts, costs, and exact-match metrics independently.

PRA supplies a neural execution substrate for a PageIndex-like index: every tree node can be a stable reference; selection may be an explicit LLM decision, a learned router, or a cascade. The most informative comparison stratifies questions by exact lookup, cross-reference, multi-section synthesis, and cross-document reasoning.

\subsection{GraphRAG and knowledge-graph retrieval}
GraphRAG-style systems extract entities and relations, form communities, generate community summaries, and answer through local or global graph search. Graph structure supports multi-hop and corpus-level thematic questions that flat vector retrieval misses. PRA references need not point only to text spans: they can address graph nodes, neighborhoods, query results, or generated summaries. The challenge is invalidation and provenance when derived graph objects change.

\subsection{PaperTrail: claim--evidence provenance}
The 2026 PaperTrail interface decomposes LLM answers and source documents into claims and evidence, exposing supported claims, unsupported assertions, and omitted source information \citep{martinboyle2026papertrail}. It is not a long-context architecture, but it contributes a missing evaluation and interaction layer. PRA's URI references and route traces could generate claim-level provenance automatically: which object, span, and layer-level memory influenced each answer segment? Attention alone is not causal proof, so this should be combined with counterfactual removal and replacement.

\section{Category VII: Agents with Search Tools}
\subsection{ReAct}
ReAct interleaves natural-language reasoning traces with actions and observations \citep{yao2023react}. For search tasks, the policy produces a query, receives snippets or pages, updates its state, and repeats. Its experiments covered knowledge-intensive QA and interactive decision environments. The insight is that reasoning can change what information should be acquired next.

PRA shares progressive disclosure but moves a restricted class of actions closer to attention. ReAct remains preferable when actions alter the world, require arbitrary APIs, or demand explicit multi-step planning. PRA is attractive for frequent, low-level, read-only reference operations. A hybrid policy should escalate from latent cache lookup to local structural navigation to external search only when confidence remains low.

\subsection{Self-Ask with Search}
Self-Ask decomposes a question into follow-up questions and delegates them to a search engine \citep{press2022selfask}. This exposes multi-hop retrieval structure and provides a clean experimental template: count subquestions, search calls, evidence hops, and final accuracy. PRA can represent follow-up targets as generated references, but must avoid treating free-form external search as if it were a deterministic URI resolution.

\subsection{IRCoT}
IRCoT interleaves retrieval with chain-of-thought so intermediate reasoning generates improved retrieval queries \citep{trivedi2023ircot}. It showed gains on multi-step QA compared with retrieve-then-read pipelines. The relation to PRA is direct: upper-layer or decoding-state representations could trigger successive memory expansions, approximating iterative retrieval without serializing every query into text. A controlled experiment compares textual query rewriting against latent query vectors with the same retrieval backend.

\subsection{Toolformer}
Toolformer self-supervises a language model to insert API calls where they improve token prediction \citep{schick2023toolformer}. It demonstrates learned triggering and result integration. PRA can borrow the training principle: sample candidate reference accesses, retain calls that reduce loss, and train a router/gate from those pseudo-labels. Toolformer still emits discrete calls and observations; PRA may execute access internally.

\subsection{WebGPT and browser agents}
WebGPT trains a model to search, click, scroll, quote sources, and answer with citations, using demonstrations and human feedback \citep{nakano2021webgpt}. It established a paper trail from browser-assisted QA to modern research agents. The key system variables are action count, source quality, citation correctness, and browsing latency. PRA cannot replace discovery over an open web; it can cache and address already discovered pages, sections, and evidence objects across subsequent reasoning or users.

\subsection{Recursive Language Models and context-as-environment}
Recursive language-model approaches treat an oversized context as an external environment over which a model writes code or issues recursive subcalls. Rather than reading everything, the model samples, searches, partitions, summarizes, and recursively invokes itself. This is the most general form of progressive context access, but also the most expensive and difficult to train. PRA can be seen as a differentiable fast path for recurring operations that an RLM or agent discovers are common.

\subsection{Agent--PRA escalation architecture}
\begin{figure}[h]
\centering
\begin{tikzpicture}[node distance=7mm, every node/.style={font=\small}, box/.style={draw,rounded corners,align=center,minimum width=35mm,minimum height=8mm}, arr/.style={-{Latex[length=2mm]}}]
\node[box,fill=green!8] (local) {local self-attention};
\node[box,fill=blue!8,below=of local] (pra) {PRA cached URI memory};
\node[box,fill=orange!10,below=of pra] (tree) {structured index traversal};
\node[box,fill=red!8,below=of tree] (agent) {external search/tool agent};
\node[right=15mm of pra,align=left] (cost) {increasing latency, cost,\\generality, and audit burden};
\draw[arr] (local)--node[right]{insufficient}(pra);
\draw[arr] (pra)--node[right]{low confidence / miss}(tree);
\draw[arr] (tree)--node[right]{source unknown}(agent);
\draw[arr,dashed] (agent.west) .. controls +(-15mm,0) and +(-15mm,0) .. node[left]{cache discoveries as references} (pra.west);
\draw[arr] (pra)--(cost);
\end{tikzpicture}
\caption{A tiered architecture uses the cheapest adequate access mechanism and caches expensive discoveries.}
\end{figure}

\fi
\clearpage
\begin{landscape}
\thispagestyle{empty}
\section{One-Page Architecture Landscape}
\begin{center}
{\LARGE \textbf{Landscape of Long-Context and Retrieval Architectures}}\\[5mm]
\begin{tikzpicture}[x=1cm,y=1cm, every node/.style={font=\scriptsize}, method/.style={draw,rounded corners,align=center,minimum width=2.7cm,minimum height=.75cm}, cat/.style={draw,very thick,rounded corners,inner sep=4mm}, arr/.style={-{Latex[length=2mm]},thick}]
\node[method,fill=gray!10] (full) at (0,6) {Dense long context\\RoPE extension};
\node[method,fill=gray!10] (sparse) at (3.2,6) {Longformer / BigBird\\Performer};
\node[method,fill=yellow!12] (txl) at (0,4) {Transformer-XL\\Compressive Transformer};
\node[method,fill=yellow!12] (rmt) at (3.2,4) {RMT / Infini-attention};
\node[method,fill=green!10] (dnc) at (6.4,6) {Memory Nets / DNC\\Memformer};
\node[method,fill=green!10] (memo) at (9.6,6) {Memorizing\\Transformer};
\node[method,fill=blue!10] (rag) at (6.4,4) {REALM / RAG\\kNN-LM};
\node[method,fill=blue!10] (retro) at (9.6,4) {RETRO\\retrieval cross-attn};
\node[method,fill=cyan!12] (unlim) at (12.8,6) {Unlimiformer\\token-state MIPS};
\node[method,fill=cyan!12] (paged) at (16,6) {PagedAttention\\KV cache systems};
\node[method,fill=purple!10] (raptor) at (12.8,4) {RAPTOR / MemWalker\\PageIndex};
\node[method,fill=purple!10] (graph) at (16,4) {GraphRAG\\PaperTrail provenance};
\node[method,fill=red!8] (react) at (3.2,1.7) {ReAct / Self-Ask\\IRCoT};
\node[method,fill=red!8] (tools) at (6.4,1.7) {Toolformer / WebGPT\\RLMs};
\node[method,fill=orange!20,very thick,minimum width=4.2cm,minimum height=1cm] (pra) at (11.2,1.7) {PRAttention\\URI naming + progressive\\latent KV materialization};
\node[method,fill=orange!8] (hybrid) at (15.5,1.7) {Proposed hybrid\\router + tree\\+ agent escalation};
\draw[arr] (full)--(sparse); \draw[arr] (txl)--(rmt); \draw[arr] (dnc)--(memo); \draw[arr] (rag)--(retro); \draw[arr] (unlim)--(pra); \draw[arr] (paged)--(pra); \draw[arr] (raptor)--(pra); \draw[arr] (graph)--(hybrid); \draw[arr] (react)--(tools); \draw[arr] (tools.south) to[bend right=16] (hybrid.south); \draw[arr] (retro)--(pra); \draw[arr] (rmt)--(pra); \draw[arr] (pra)--(hybrid);
\node[align=center] at (9,-.1) {\textbf{Main axes:} eager $\leftrightarrow$ demand-driven access; flat tokens $\leftrightarrow$ named structure; text $\leftrightarrow$ latent KV;\\static one-shot selection $\leftrightarrow$ progressive reasoning; transient context $\leftrightarrow$ reusable cache; opaque use $\leftrightarrow$ causal provenance.};
\end{tikzpicture}
\end{center}
\end{landscape}

\section{Cross-Architecture Comparison}
\small
\begin{longtable}{>{\raggedright\arraybackslash}p{2.6cm}>{\raggedright\arraybackslash}p{2.5cm}>{\raggedright\arraybackslash}p{2.3cm}>{\raggedright\arraybackslash}p{2.3cm}>{\raggedright\arraybackslash}p{2.4cm}>{\raggedright\arraybackslash}p{2.5cm}}
\toprule
\textbf{Family} & \textbf{Store} & \textbf{Selection time} & \textbf{Form} & \textbf{Main strength} & \textbf{Characteristic failure}\\\midrule
Dense long context & prompt tokens & before forward pass & token KVs & general interactions & quadratic cost, distraction\\
Sparse attention & prompt tokens & architecture or task mask & sparse token KVs & long visible sequence & missing sparse paths\\
Recurrence & prior segments & chronological & hidden/KV state & streaming continuity & lossy or indiscriminate history\\
Differentiable memory & learned slots/episodes & per layer/step & latent vectors & trainable read/write & optimization and identity\\
RAG/REALM & document corpus & before or per-token generation & retrieved text encodings & fresh external knowledge & chunking and retrieval miss\\
RETRO & token chunks & chunk boundary & cross-attention encodings & scale via nonparametric data & index and corpus dependence\\
KV retrieval & hidden/KV index & attention time & selected states & efficient latent access & provenance and stale cache\\
Tree/graph index & structural nodes & iterative or one-shot & summaries, spans, or subgraphs & multigranular reasoning & index construction errors\\
Search agent & open tools/web & iterative decisions & observations or text & maximum flexibility & latency, cost, tool errors\\
PRA & URI objects and KV caches & layer/step, demand-driven & selected per-layer KV & identity, reuse, progressive access & router training and runtime complexity\\
\bottomrule
\end{longtable}
\normalsize

\section{Adapting Original Experiments to PRA}
\subsection{A benchmark matrix}
No single benchmark distinguishes all mechanisms. A credible program combines:
\begin{enumerate}
\item \textbf{Language preservation}: WikiText-2/103 or C4 slices before and after memory adaptation.
\item \textbf{Exact long-range retrieval}: passkeys, variable needles, UUID/value lookup, and repeated distractors.
\item \textbf{Natural long-document understanding}: NarrativeQA, Qasper, QuALITY, GovReport, arXiv/PubMed summarization, and LongBench-style tasks.
\item \textbf{Open-domain retrieval}: Natural Questions, TriviaQA, HotpotQA, 2WikiMultiHopQA, MuSiQue.
\item \textbf{Structured documents}: FinanceBench, legal filings, technical manuals, code repositories, and table/figure references.
\item \textbf{Agentic discovery}: web or corpus search tasks where the correct source is initially unknown.
\end{enumerate}

\subsection{Common conditions}
For each task, compare: local-only SelfAttention; longer-context SelfAttention; sparse attention; one-shot RAG; iterative RAG; latent token retrieval; PRA with explicit handles; PRA with learned discovery; structural-tree retrieval; and agentic search. Hold corpus, tokenizer, generation model, evidence budget, and evaluation prompts fixed where meaningful.

\subsection{PRA-specific causal matrix}
\begin{center}
\begin{tabular}{lll}
\toprule
Condition & What it tests & Expected interpretation\\\midrule
Valid & complete system & task utility\\
Disabled branch & any memory contribution & $\Delta_{use}$\\
Shuffled content & semantic dependence & $\Delta_{content}$\\
Irrelevant matched length & distractor resistance & router specificity\\
Random URI, same content & identity dependence & naming shortcut\\
Same URI, wrong version & cache invalidation & staleness safety\\
Oracle object & selection upper bound & router headroom\\
All evidence local & access overhead & efficiency frontier\\
Agent oracle search & discovery upper bound & unresolved-source headroom\\
\bottomrule
\end{tabular}
\end{center}

\subsection{Metrics beyond task accuracy}
Report task quality, retrieval recall, ranking, reference causality, preservation, and systems cost together. Systems measurements must separate resolve time, encode time, cache lookup, routing, attention, generation, and synchronization. Cache hit rate without latency reduction is not useful; latency reduction without answer preservation is not success.

\section{Research Lineage and Influence}
The field can be read as a sequence of boundary shifts. Memory Networks moved relevant facts outside the recurrent state. Transformer-XL moved context across segments. REALM and RAG moved knowledge outside parameters. RETRO and Memorizing Transformer moved retrieved information closer to hidden-state computation. Unlimiformer recast cross-attention as nearest-neighbor search. PagedAttention made KV storage a virtual-memory systems problem. RAPTOR, MemWalker, and PageIndex restored document hierarchy. ReAct, IRCoT, Toolformer, and WebGPT made information acquisition an action policy. PRA's proposed step is to make \emph{named reference resolution and reusable latent materialization} a first-class attention capability.

This lineage also tempers novelty claims. Cross-attention over external KV is not new; retrieval during model computation is not new; sparse top-$k$ selection is not new; and caches are not new. PRA's distinct contribution is the package and interface: URI-level identity, explicit reference syntax, separation of naming and materialization, per-layer cache objects, progressive triggering, and a causal evaluation doctrine. Its scientific value depends on showing that this package creates measurable advantages unavailable from simpler combinations.

\section{Design Extensions Derived from the Survey}
\subsection{Hierarchical Progressive Retrieval Attention}
Represent a document as a URI tree. Each node stores a summary KV and child handles. The router first attends to coarse nodes, then expands children only when posterior entropy or expected utility warrants it. This unifies RAPTOR/PageIndex with PRA while preserving traceability.

\subsection{Agent-distilled routing}
Use a capable search agent to generate trajectories: queries, selected sources, section paths, and evidence spans. Convert these into reference-routing supervision. The deployed PRA model learns a low-latency approximation; the agent remains an escalation path for out-of-index or ambiguous cases.

\subsection{Multi-fidelity object caches}
For each URI maintain tiers: metadata, summary vector, compressed KV, full KV, and raw source. A learned controller chooses the cheapest tier expected to reduce loss. This connects Compressive Transformer, Infini-attention, PagedAttention, and PRA.

\subsection{Reference compilation}
Repeated agent/tool operations can be compiled into stable references. A SQL result, code symbol graph, dashboard panel, or document query becomes a versioned object with dependency metadata. PRA then retrieves the result cheaply until invalidated. This is especially relevant to data-analytics agents.

\subsection{Causal provenance graph}
Record reference selection, cache version, attended spans, generated claims, and counterfactual effects. Join these into a graph compatible with PaperTrail-like claim-evidence views. The strongest evidence edge is not raw attention weight but measured output change under controlled removal or substitution.

\section{Discussion}
\subsection{The real design choice is the access contract}
Debates framed as ``long context versus RAG'' hide the more useful question: what contract exists between computation and external information? Dense context offers anonymous token positions. RAG offers retrieved passages. agents offer observations from named actions. PRA proposes persistent objects with identity, version, representation, and cache lifecycle. Such contracts determine not only quality but governance, sharing, invalidation, debugging, and cost.

\subsection{When PRA should not win}
PRA is unlikely to dominate when the entire source is short, every token is densely relevant, hardware kernels make dense attention cheap, or the source is unknown and requires open-ended discovery. It also adds complexity where a simple one-shot retriever is adequate. Negative results in these regimes would clarify the architecture rather than invalidate it.

\subsection{When PRA has a plausible advantage}
The strongest regime has large, repeatedly used, naturally addressable objects; sparse relevance; many queries over shared sources; stable or versioned content; and tasks where relevance becomes clearer only after intermediate processing. Examples include technical documentation, analytics catalogs, code repositories, longitudinal records, enterprise knowledge, and embodied-agent maps.

\subsection{Training difficulty}
A new memory branch can be ignored, overused, or used as a generic adapter. Near-zero gate initialization, frozen-base staged training, auxiliary selection losses, corrupted-reference negatives, and curriculum by distance are sensible. However, explicit supervision may bias the model toward the indexer's notion of relevance. End-to-end loss and causal ablations remain necessary.

\subsection{Object identity versus semantic similarity}
Semantic retrieval answers ``what looks relevant?'' Explicit addressing answers ``which object is being referred to?'' Real systems need both. An explicit handle can restrict search to an object namespace; semantic routing can select parts within it; an agent can discover a missing object. Treating one mechanism as universally sufficient is a category error.

\subsection{Impact on agents}
PRA does not eliminate the agent loop. It can internalize routine read-only actions, reduce repeated tokenization and tool latency, and provide a memory substrate for agent discoveries. The agent retains planning, source discovery, side effects, and exception handling. The long-term architecture is likely layered rather than exclusive.

\section{Conclusion}
Long-context modeling is not one problem. It is a family of access problems spanning compute graphs, state persistence, external memory, retrieval, indexing, runtime allocation, provenance, and action. The seven categories reviewed here form a continuum from eager anonymous tokens to progressive named actions. PRA occupies a strategically interesting middle point: more structured and reusable than flat long context or one-shot RAG, but potentially cheaper and more differentiable than general search agents.

The current evidence establishes implementation feasibility, ordinary-language compatibility under approximate parameter matching, and preservation under a frozen reference-path adaptation strategy. It does not establish content-selective reference use: the fixed three-seed follow-up leaves the selector near chance, valid references barely outperform shuffled or irrelevant controls, and oracle content does not materially improve loss. The most valuable next experiments are therefore a trainable router and reference-required natural tasks, followed by five-seed replication, hierarchical and agentic baselines, and synchronized measurements of latency, memory, encoding, and cache reuse. If those experiments succeed, PRA could become a useful abstraction for compiling repeated referential work into a model-native memory interface. If they fail, the same experimental program will still clarify where dense context, RAG, structural navigation, and agentic search should meet.

\section*{Online Resources}
PRA source manuscripts used for this tutorial were supplied by the author. Live links for external works are embedded in the bibliography through DOI or arXiv URLs. The PageIndex discussion refers to its public repository and documentation; claims should be independently reproduced.

\bibliographystyle{plainnat}
\bibliography{references}
\end{document}
